%%\documentclass[a4paper,11pt]{article}
\documentclass[11pt]{article}
\usepackage[left=2.5cm, right=2cm, top=1.5cm, bottom=3cm]{geometry}
\usepackage[utf8]{inputenc}
\usepackage{titlesec}
\usepackage{enumitem}
\usepackage{graphicx}
\usepackage{amsmath}
\usepackage{amsfonts}
\usepackage{amssymb}
\usepackage{tabularx}
\usepackage{array}
\usepackage{longtable}
\usepackage{float}
\usepackage{fancyhdr}
\pagestyle{fancy}
\fancyhf{}
\setlength{\headheight}{35pt}
\setlength{\headsep}{35pt}
\fancyhead[L]{\includegraphics[width=18cm]{picture/logo.pdf}}
\renewcommand{\headrulewidth}{0pt}
\fancyfoot[C]{\thepage}
\setlength{\footskip}{5pt}
\title{Informe del Proyecto REGALÍA}

% Let long monospace identifiers (endpoints, error codes) wrap inside narrow cells.
\let\oldtexttt\texttt
\renewcommand{\texttt}[1]{\oldtexttt{\hyphenchar\font=45\relax #1}}
\renewcommand{\_}{\textunderscore\allowbreak}
\setlength{\emergencystretch}{3em}
\renewcommand{\figurename}{Figura}

\newcolumntype{L}{>{\raggedright\arraybackslash}X}
\newcolumntype{P}[1]{>{\raggedright\arraybackslash}p{#1}}
\renewcommand{\arraystretch}{1.3}

% ---------------------------------------------------------------------------
% Macros for the test design and execution tables.
% Design row:    \caso{ID}{Case and requirement}{Inputs}{Actions}{Expected}{Acceptance}
% Execution row: \fila{ID}{Obtained result}{Status}{Evidence}{Defect ID}
% Screenshot:    \captura{ID}{Description} -> loads evidencias/ID.png if it exists
% ---------------------------------------------------------------------------
\newcommand{\caso}[6]{\oldtexttt{#1} & #2 & #3 & #4 & #5 & #6 \\ \hline}
\newenvironment{disenio}{%
  \sloppy\scriptsize\setlength{\tabcolsep}{3pt}
  \begin{longtable}{|P{1.8cm}|P{2.2cm}|P{2.75cm}|P{2.85cm}|P{2.8cm}|P{2.65cm}|}
  \hline
  \textbf{ID} & \textbf{Caso / requisito} & \textbf{Entradas} & \textbf{Acciones} & \textbf{Resultado esperado} & \textbf{Criterio de aceptación} \\ \hline
  \endhead
}{\end{longtable}\fussy\normalsize}

\newcommand{\fila}[5]{\oldtexttt{#1} & #2 & #3 & #4 & #5 \\ \hline}
\newenvironment{ejecucion}{%
  \sloppy\footnotesize\setlength{\tabcolsep}{3pt}
  \begin{longtable}{|P{2.05cm}|P{6.3cm}|P{1.6cm}|P{3.4cm}|P{1.5cm}|}
  \hline
  \textbf{ID} & \textbf{Resultado obtenido} & \textbf{Estado} & \textbf{Evidencia} & \textbf{Defecto} \\ \hline
  \endhead
}{\end{longtable}\fussy\normalsize}

% Ficha individual numerada (2.7.4.1 .. 2.7.4.107) para Ejecución de pruebas:
% \casoejec{ID}{Título}{Resultado obtenido}{Estado}{Evidencia / Log}{Defecto}
\newcommand{\casoejec}[6]{%
  \paragraph{\oldtexttt{#1} --- #2}\nopagebreak[4]%
  \par\noindent
  \begin{minipage}{\textwidth}
    \sloppy
    {\footnotesize
      \textbf{Estado:} \if\relax\detokenize{#4}\relax\underline{\hspace{2.5cm}}\else\textbf{#4}\fi
      \hfill
      \textbf{Defecto:} \if\relax\detokenize{#6}\relax\underline{\hspace{2.5cm}}\else #6\fi\par
      \vspace{4pt}
      \textbf{Resultado obtenido:}
      \if\relax\detokenize{#3}\relax
        \par\vspace{8pt}\noindent\rule{\textwidth}{0.3pt}\par\vspace{8pt}\noindent\rule{\textwidth}{0.3pt}\par
      \else
        #3\par
      \fi
      \vspace{4pt}
      \textbf{Evidencia (log / reporte):} #5\par
      \vspace{4pt}
      \IfFileExists{evidencias/#1.png}%
        {\begin{center}\includegraphics[width=0.78\textwidth,height=0.25\textheight,keepaspectratio]{evidencias/#1.png}\end{center}}%
        {\noindent\fbox{\parbox[c][1.35cm][c]{\dimexpr\textwidth-2\fboxsep-2\fboxrule\relax}{\centering\scriptsize Evidencia gr\'afica: \texttt{evidencias/#1.png}}}}%
    }
    \vspace{4pt}
    \noindent\rule{\textwidth}{0.4pt}\par
    \fussy
  \end{minipage}\par
  \vspace{4pt}
}

% GUÍA EQUIPO: guardar cada captura como evidencias/CP-XXX.png; \captura la inserta automáticamente si el archivo existe.
\newcommand{\captura}[2]{%
  \begin{figure}[H]
    \centering
    \IfFileExists{evidencias/#1.png}%
      {\includegraphics[width=0.8\textwidth,height=0.32\textheight,keepaspectratio]{evidencias/#1.png}}%
      {\fbox{\parbox[c][2.5cm][c]{0.8\textwidth}{\centering\small Evidencia gr\'afica #1 (\texttt{evidencias/#1.png})}}}
    \caption{#1: #2}
  \end{figure}}

\begin{document}

\titleformat{\section}{\normalfont\Large\bfseries}{\thesection}{1em}{}
\renewcommand{\thesection}{\Roman{section}}
\begin{center}
    \textbf{\underline{\normalsize INFORME DE GUÍA PRÁCTICA}}
\end{center}

\section{PORTADA}

\begin{tabbing}
\textbf{Tema:}\hspace{7cm} \= APE GPIS 4: Diseño y ejecución de pruebas\\
\textbf{Carrera:}\hspace{7cm} \= SOFTWARE \\
\textbf{Unidad de Organización Curricular:} \hspace{1cm} \= \ PROFESIONAL\\
\textbf{Nivel y Paralelo:} \> 6 - A\\
\textbf{Alumnos participantes:}  \> Auz Mayorga José Mateo\\
 \> Chicaiza Llerena Kerly Mishelle\\
 \> Jimenez Chango Saimol Steph\\
 \> Ramos Jacome Carlos Giovanni\\
\textbf{Asignatura:} \> Gestión de pruebas e integración de software \\
\textbf{Docente:} \> Mgtr. Ing. Leonardo David Torres Valverde
\end{tabbing}

\section{INFORME DE GUÍA PRÁCTICA}
\titleformat{\section}{\normalfont\Large\bfseries}{\thesection}{1em}{}
\renewcommand{\thesection}{\arabic{section}}

\subsection{Objetivos}
\textbf{General:}\\
Diseñar y ejecutar casos de prueba que permitan validar el correcto funcionamiento de un sistema de software.

\textbf{Específicos:}\\
\begin{itemize}
    \item Analizar los requisitos funcionales (\texttt{NR-01} a \texttt{NR-29}) y escenarios de
    prueba (\texttt{NP-01} a \texttt{NP-19}) definidos en la documentación del sistema REGALÍA.
    \item Diseñar y ejecutar 107 casos de prueba individuales (\texttt{CP-001} a \texttt{CP-107}),
    abarcando variantes funcionales, de seguridad, concurrencia y rendimiento (\texttt{100/500/2000}
    pedidos), con entradas, acciones, resultados esperados y criterios de aceptación.
    \item Configurar el entorno aislado con datos sintéticos, registrando los resultados obtenidos y
    los defectos detectados con su nivel de severidad.
\end{itemize}

\subsection{Modalidad}
\begin{tabbing}
    \hspace{1cm} \= Presencial\\
\end{tabbing}

\subsection{Tiempo de duración}
\begin{tabbing}
    \hspace{1cm} \= \textbf{Presenciales:} \= 8 \\
    \hspace{1cm} \= \textbf{No presenciales:} \= 0
\end{tabbing}

\subsection{Instrucciones}
En base a los requerimientos definidos para el desarrollo del sistema realice lo siguiente:
\begin{itemize}
    \item \textbf{Análisis de requisitos:} revisar la documentación del sistema para identificar qué funcionalidades deben probarse.
    \item \textbf{Diseño de pruebas:} elaborar los casos de prueba especificando entradas, acciones, resultados esperados y criterios de aceptación.
    \item \textbf{Preparación del entorno:} asegúrese de contar con los datos de prueba y el ambiente configurado.
    \item \textbf{Ejecución de pruebas:} realizar las pruebas según lo planificado, registrando los resultados obtenidos.
    \item \textbf{Registro de incidentes:} documentar cualquier error o comportamiento inesperado en un reporte de defectos.
    \item \textbf{Verificación de resultados:} comparar los resultados reales con los esperados y concluir si las pruebas fueron exitosas.
    \item \textbf{Entrega de evidencias:} generar un documento en formato pdf en base al formato indicado.
\end{itemize}

\subsection{Listado de equipos, materiales y recursos}
Listado de equipos y materiales generales empleados en la guía práctica:\\
\begin{itemize}
    \item Computador o laptop, y un dispositivo o emulador Android para la aplicación móvil.
    \item Internet.
    \item Apuntes de clase.
    \item Diseñar y ejecutar casos de prueba que permitan validar el correcto funcionamiento de un sistema de software.
\end{itemize}
\textbf{TAC (Tecnologías para el Aprendizaje y Conocimiento) empleados en la guía práctica:}\\
$\boxtimes$ Plataformas educativas \\
$\boxtimes$ Simuladores y laboratorios virtuales \\
$\square$ Aplicaciones educativas \\
$\square$ Recursos audiovisuales \\
$\square$ Gamificación \\
$\boxtimes$ Inteligencia Artificial \\
Otros (Especifique): \underline{} \\

\subsection{Actividades por desarrollar}
Desarrollar la guía APE en base al formato y las instrucciones brindadas.

\subsection{Resultados obtenidos}
Se diseñaron y registraron resultados para \textbf{107 casos de prueba individuales}
(\texttt{CP-001} a \texttt{CP-107}) del sistema REGALÍA, integrado por una API NestJS, una
aplicación web React y una aplicación móvil Flutter. Las fichas documentan los resultados de la
última ejecución registrada para cada caso y conservan los antecedentes de las reejecuciones:

% GUÍA PARA EL EQUIPO (asignación interna y cronograma de registro; no visible en el PDF):
% - Backend (Saimol Jimenez, 27 casos: CP-001..CP-027): levantar entorno común (PostgreSQL, seed aislado), API/settings y llenar secciones/backend/ejecucion.tex + capturas.tex.
% - Frontend Web (Kerly Chicaiza, 27 casos: CP-028..CP-054): pruebas Shop web, cobertura poligonal, panel de rutas y llenar secciones/frontend/ejecucion.tex + capturas.tex.
% - Móvil (Carlos Ramos, 27 casos: CP-055..CP-081): emulador/app Flutter, entregas, fallos, retorno y llenar secciones/movil/ejecucion.tex + capturas.tex.
% - Calidad (José Mateo Auz, 26 casos: CP-082..CP-107): concurrencia, emergencias, pérdidas, carga 100/500/2000, llenar secciones/calidad/*, consolidar defectos.tex, Verificación y Conclusiones.
% - Cronograma sugerido (8h): H1 entorno y datos sintéticos; H2-6 ejecución 107 casos y capturas; H7 defectos y verificación; H8 PDF final.

\begin{center}
\begin{tabularx}{\textwidth}{|l|c|c|L|}
\hline
\textbf{Módulo del sistema} & \textbf{Rango de casos} & \textbf{Cantidad} & \textbf{Alcance evaluado} \\ \hline
Backend y reglas de negocio & \texttt{CP-001}--\texttt{CP-027} & 27 & API REST, comercio, cobros simulados, configuraciones e idempotencia. \\ \hline
Aplicación Web REGALÍA & \texttt{CP-028}--\texttt{CP-054} & 27 & Catálogo Shop, cobertura poligonal, horarios, historial y panel de rutas. \\ \hline
Aplicación Móvil & \texttt{CP-055}--\texttt{CP-081} & 27 & Seguridad RBAC, asignación de entregas, fallos, retiros y retorno. \\ \hline
Calidad y Rendimiento & \texttt{CP-082}--\texttt{CP-107} & 26 & Agrupación de rutas, emergencias, pérdidas, concurrencia y carga \texttt{100/500/2000}. \\ \hline
\textbf{Total} & \texttt{CP-001}--\texttt{CP-107} & \textbf{107} & Cobertura integral del sistema REGALÍA. \\ \hline
\end{tabularx}
\end{center}

% ===========================================================================
\subsubsection{Análisis de requisitos}
Se analizaron los 29 requisitos del sistema (\texttt{NR-01} a \texttt{NR-29}) y los 19 escenarios
principales (\texttt{NP-01} a \texttt{NP-19}) correspondientes a los módulos comercial, geográfico,
de planificación de rutas, operación móvil y calidad. La trazabilidad entre requisitos y casos de
prueba es la siguiente:

\begin{center}
\footnotesize
\begin{tabularx}{\textwidth}{|P{4.2cm}|L|}
\hline
\textbf{Grupo de requisitos (\texttt{NR}/\texttt{NP})} & \textbf{Casos de prueba asignados} \\ \hline
\texttt{NR-01}--\texttt{NR-04} Compras ADMIN, invitado, pago externo y contacto & \texttt{CP-001}--\texttt{CP-007}, \texttt{CP-010}--\texttt{CP-012} \\ \hline
\texttt{NR-05}--\texttt{NR-09} Precios/IVA congelados, historial único y sin recibo & \texttt{CP-008}--\texttt{CP-009}, \texttt{CP-028}--\texttt{CP-030} \\ \hline
\texttt{NR-10} Operación Shop, idempotencia y pedidos sin fecha programada & \texttt{CP-022}--\texttt{CP-026}, \texttt{CP-052}--\texttt{CP-054} \\ \hline
\texttt{NR-11}--\texttt{NR-15} Cobertura de tienda, horarios y trazado & \texttt{CP-013}--\texttt{CP-016}, \texttt{CP-031}--\texttt{CP-041}, \texttt{CP-049}--\texttt{CP-051} \\ \hline
\texttt{NR-16}--\texttt{NR-22} Planificación diaria, margen horario, urgencias y motor vial & \texttt{CP-017}--\texttt{CP-021}, \texttt{CP-027}, \texttt{CP-042}--\texttt{CP-048}, \texttt{CP-061}--\texttt{CP-066}, \texttt{CP-082}--\texttt{CP-085} \\ \hline
\texttt{NR-23} Seguridad y control de acceso por recurso & \texttt{CP-055}--\texttt{CP-060} \\ \hline
\texttt{NR-24}--\texttt{NR-28} Fallos, retiro, cierre/retorno, emergencias y pérdidas & \texttt{CP-067}--\texttt{CP-081}, \texttt{CP-086}--\texttt{CP-099} \\ \hline
\texttt{NR-29} y no funcionales: concurrencia, rollback y rendimiento \texttt{100/500/2000} & \texttt{CP-100}--\texttt{CP-107} \\ \hline
\end{tabularx}
\end{center}

\newpage
% ===========================================================================
\subsubsection{Diseño de pruebas}
Cada caso del catálogo incluye entradas, acciones, resultado esperado y criterio de aceptación
verificable. Se utilizan cuentas sintéticas aisladas (\texttt{ADMIN}, \texttt{CLIENTE-A},
\texttt{CLIENTE-B}, \texttt{REPARTIDOR-A}, \texttt{REPARTIDOR-B} e invitado) y las tres tarjetas
ficticias del contrato del Shop con \texttt{SIMULATED\_PAYMENTS\_ENABLED=true}:

\begin{center}
\footnotesize
\begin{tabularx}{\textwidth}{|l|l|l|c|c|L|}
\hline
\textbf{Fixture} & \textbf{Número ficticio} & \textbf{Titular} & \textbf{Vence} & \textbf{CVV} & \textbf{Saldo inicial} \\ \hline
DEMO\_APPROVED & 9000 0000 0000 0001 & PRUEBA APROBADA & 12/30 & 111 & 50000 centavos (\$500,00) \\ \hline
DEMO\_LIMITED & 9000 0000 0000 0002 & PRUEBA LIMITADA & 12/30 & 222 & 2000 centavos (\$20,00) \\ \hline
DEMO\_EMPTY & 9000 0000 0000 0003 & PRUEBA SIN SALDO & 12/30 & 333 & 0 centavos \\ \hline
\end{tabularx}
\end{center}

Las credenciales de prueba no se almacenan en este documento. Para el ejemplo monetario estándar se
emplea subtotal \texttt{7300} centavos, IVA simulado \texttt{15\,\%} (\texttt{1095}), envío \texttt{0}
y total \texttt{8395}.

\clearpage
% BACKEND - diseño de 27 casos (CP-001..CP-027)
% GUÍA EQUIPO: Responsable asignado: Saimol Jimenez (Backend).
\paragraph{Módulo Backend y reglas de negocio (27 casos: CP-001 a CP-027).}

\begin{disenio}
\caso{CP-001}{\textbf{Compra ADMIN para cliente con destinatario de regalo.} NR-01, NR-02; NP-01 (C01.1, Alta)}{Token ADMIN; CLIENT-A con contactPhone distinto al de giftRecipient; ítem en stock; pago CASH; Idempotency-Key UUIDv4.}{1) Enviar POST /\allowbreak{}api/\allowbreak{}v1/\allowbreak{}orders/\allowbreak{}admin con buyer REGISTERED\_CLIENT y giftRecipient. 2) Consultar GET /\allowbreak{}api/\allowbreak{}v1/\allowbreak{}orders/\allowbreak{}:\allowbreak{}id y /\allowbreak{}receipt (ADMIN). 3) Consultar GET /\allowbreak{}api/\allowbreak{}v1/\allowbreak{}shop/\allowbreak{}orders/\allowbreak{}:\allowbreak{}id (CLIENT-A).}{HTTP 201; pedido conserva comprador y destinatario separados; ADMIN lee pedido/\allowbreak{}comprobante y CLIENT-A consulta su compra.}{Titularidad asociada a CLIENT-A, datos de giftRecipient independientes, stock descontado y comprobante comercial coherente.}
\caso{CP-002}{\textbf{Rechazo compra ADMIN cliente sin teléfono de comprador.} NR-01, NR-02; NP-01 (C01.2, Alta)}{Token ADMIN; payload POST /\allowbreak{}api/\allowbreak{}v1/\allowbreak{}orders/\allowbreak{}admin con buyer.\allowbreak{}type=\allowbreak{}REGISTERED\_CLIENT omitiendo contactPhone; Idempotency-Key UUIDv4.}{1) Registrar conteo inicial de pedidos, comprobantes y stock del producto. 2) Enviar POST /\allowbreak{}api/\allowbreak{}v1/\allowbreak{}orders/\allowbreak{}admin sin contactPhone del comprador. 3) Consultar conteos y stock posterior.}{HTTP 400 de validación rechazando la creación del pedido por falta de contactPhone obligatorio del comprador.}{Sin pedido creado, sin PurchaseReceipt ni descuento de stock; conteos inicial y final idénticos.}
\caso{CP-003}{\textbf{Compra ADMIN invitado y lectura de comprobante.} NR-01, NR-02, NR-06, NR-09; NP-02 (C02.1, Alta)}{Token ADMIN; buyer.\allowbreak{}type=\allowbreak{}GUEST con name y contactPhone válidos; giftRecipient opcional; pago CASH; Idempotency-Key UUIDv4.}{1) Enviar POST /\allowbreak{}api/\allowbreak{}v1/\allowbreak{}orders/\allowbreak{}admin como invitado. 2) Invocar GET /\allowbreak{}api/\allowbreak{}v1/\allowbreak{}orders/\allowbreak{}:\allowbreak{}id/\allowbreak{}receipt con token ADMIN. 3) Verificar estructura imprimible del comprobante.}{HTTP 201 en creación; GET /\allowbreak{}receipt retorna 200 con número, ítems, importes en centavos, comprador invitado y aviso no fiscal.}{Comprobante corresponde exactamente a la venta y contiene los campos necesarios para impresión ADMIN sin exigir cuenta de usuario.}
\caso{CP-004}{\textbf{Rechazo compra ADMIN invitado sin teléfono obligatorio.} NR-01, NR-02, NR-06, NR-09; NP-02 (C02.2, Alta)}{Token ADMIN; buyer.\allowbreak{}type=\allowbreak{}GUEST con name válido y contactPhone omitido o vacío; Idempotency-Key UUIDv4.}{1) Verificar stock y total de pedidos/\allowbreak{}usuarios antes del envío. 2) Enviar POST /\allowbreak{}api/\allowbreak{}v1/\allowbreak{}orders/\allowbreak{}admin sin contactPhone del invitado. 3) Reconsultar stock y tablas.}{HTTP 400 rechazando el pedido de invitado sin teléfono obligatorio del comprador.}{No se persiste pedido ni comprobante, no se crea usuario ni se descuenta stock del producto.}
\caso{CP-005}{\textbf{Compra invitado con teléfono de cliente sin vínculo.} NR-01, NR-02, NR-06, NR-09; NP-02 (C02.3, Alta)}{Token ADMIN; buyer.\allowbreak{}type=\allowbreak{}GUEST cuyo contactPhone coincide exactamente con el teléfono de CLIENT-A; Idempotency-Key UUIDv4.}{1) Crear venta ADMIN de invitado con teléfono de CLIENT-A. 2) Consultar GET /\allowbreak{}api/\allowbreak{}v1/\allowbreak{}shop/\allowbreak{}orders y /\allowbreak{}shop/\allowbreak{}orders/\allowbreak{}:\allowbreak{}id con JWT de CLIENT-A. 3) Verificar cuentas en BD.}{Venta creada con userId=\allowbreak{}null; CLIENT-A no ve el pedido en su listado y su consulta directa devuelve 404/\allowbreak{}403.}{No se crea vínculo por coincidencia telefónica ni cuenta nueva; aislamiento de acceso preservado.}
\caso{CP-006}{\textbf{Pago externo declarado ADMIN en efectivo.} NR-04, NR-06; NP-02, NP-16 (C03.1, Alta)}{Token ADMIN; pedido válido con payment=\allowbreak{}\{method:\allowbreak{}'CASH', note:\allowbreak{}'Pago en caja'\}; Idempotency-Key UUIDv4; sin datos de tarjeta.}{1) Enviar POST /\allowbreak{}api/\allowbreak{}v1/\allowbreak{}orders/\allowbreak{}admin con payment.\allowbreak{}method=\allowbreak{}CASH. 2) Consultar respuesta, GET /\allowbreak{}api/\allowbreak{}v1/\allowbreak{}orders/\allowbreak{}:\allowbreak{}id/\allowbreak{}receipt y registro de pago declarado.}{HTTP 201 con kind=\allowbreak{}ADMIN\_DECLARED\_EXTERNAL, status=\allowbreak{}DECLARED, method=\allowbreak{}CASH, declaredById y declaredAt; sin débito simulado.}{Declaración y comprobante coherentes en CASH; cero campos bancarios o cobros de tarjeta simulada.}
\caso{CP-007}{\textbf{Pago externo declarado ADMIN por transferencia.} NR-04, NR-06; NP-02, NP-16 (C03.2, Alta)}{Token ADMIN; pedido válido con payment=\allowbreak{}\{method:\allowbreak{}'TRANSFER'\} sin referencia bancaria ni comprobante bancario; Idempotency-Key UUIDv4.}{1) Enviar POST /\allowbreak{}api/\allowbreak{}v1/\allowbreak{}orders/\allowbreak{}admin con payment.\allowbreak{}method=\allowbreak{}TRANSFER sin referencia bancaria. 2) Consultar GET /\allowbreak{}api/\allowbreak{}v1/\allowbreak{}orders/\allowbreak{}:\allowbreak{}id/\allowbreak{}receipt.}{HTTP 201 sin exigir referencia bancaria; status=\allowbreak{}DECLARED, method=\allowbreak{}TRANSFER, declaredById y declaredAt; comprobante sin verificación bancaria.}{Se persiste transferencia declarada por ADMIN sin presentar validación bancaria automática ni datos bancarios inventados.}
\caso{CP-008}{\textbf{Cambio de tasa IVA tras cotizar (409 SHOP\_SETTINGS\_CHANGED).} NR-05, NR-06; NP-01, NP-02, NP-03 (C04.1, Alta)}{Token CLIENT-A; carrito válido; rateBps inicial 1500 (15\%); PATCH admin/\allowbreak{}shop-settings cambia rateBps a 1200; Idempotency-Key UUIDv4.}{1) Ejecutar POST /\allowbreak{}api/\allowbreak{}v1/\allowbreak{}shop/\allowbreak{}checkout/\allowbreak{}quote como CLIENT-A. 2) Cambiar tax.\allowbreak{}rateBps vía PATCH /\allowbreak{}api/\allowbreak{}v1/\allowbreak{}admin/\allowbreak{}shop-settings. 3) Enviar POST /\allowbreak{}api/\allowbreak{}v1/\allowbreak{}shop/\allowbreak{}checkout con el quoteToken previo.}{POST /\allowbreak{}api/\allowbreak{}v1/\allowbreak{}shop/\allowbreak{}checkout responde HTTP 409 con error.\allowbreak{}code=\allowbreak{}SHOP\_SETTINGS\_CHANGED rechazando condiciones obsoletas.}{Sin pedido creado, sin PurchaseReceipt, sin descuento de stock ni débito simulado en la cuenta del cliente.}
\caso{CP-009}{\textbf{Cambio de precio tras cotizar e inmutabilidad 15\% IVA.} NR-05, NR-06; NP-01, NP-02, NP-03 (C04.2, Alta)}{Ítems por 7300 centavos, IVA 15\% (1095), total 8395; cambio de precio tras cotizar y cambio de tasa/\allowbreak{}precio post-compra.}{1) Cotizar, alterar precio y confirmar \ensuremath{\rightarrow} 409 PRICE\_CHANGED. 2) Recotizar subtotal 7300 y confirmar checkout (201). 3) Cambiar tasa/\allowbreak{}precio y releer GET /\allowbreak{}api/\allowbreak{}v1/\allowbreak{}shop/\allowbreak{}orders/\allowbreak{}:\allowbreak{}id/\allowbreak{}receipt.}{Primer intento retorna 409 PRICE\_CHANGED sin efectos; compra válida fija subtotal 7300, IVA 1095, total 8395 y aviso no fiscal.}{Relecturas del comprobante tras cambios posteriores muestran importes idénticos (7300/\allowbreak{}1095/\allowbreak{}8395) y sin etiqueta de factura fiscal.}
\caso{CP-010}{\textbf{Corrección autorizada de contacto ADMIN y auditoría.} NR-03, NR-07; NP-04 (C07.1, Alta)}{Token ADMIN; orderId existente; payload \{contactRevision:\allowbreak{}0, buyerPhone:\allowbreak{}'+593990001112', recipientPhone:\allowbreak{}'+593990001113'\}.}{1) Enviar PATCH /\allowbreak{}api/\allowbreak{}v1/\allowbreak{}orders/\allowbreak{}:\allowbreak{}id/\allowbreak{}contact con contactRevision=\allowbreak{}0. 2) Consultar GET /\allowbreak{}api/\allowbreak{}v1/\allowbreak{}orders/\allowbreak{}:\allowbreak{}id/\allowbreak{}contact-history.}{HTTP 200 actualiza contacto y eleva contactRevision a 1; contact-history lista cambios con field, previousValue, nextValue, actorId y createdAt.}{Corrección operativa aplicada y auditoría append-only trazable al actor ADMIN con valores anterior y nuevo.}
\caso{CP-011}{\textbf{Aislamiento de corrección de contacto en recibo y pedido.} NR-03, NR-07; NP-04 (C07.2, Alta)}{Token ADMIN; orderId P1 (contactRevision=\allowbreak{}0) con nuevo recipientPhone; orderId P2 de control sin cambios.}{1) Capturar snapshot de GET /\allowbreak{}orders/\allowbreak{}:\allowbreak{}id1/\allowbreak{}receipt y GET /\allowbreak{}orders/\allowbreak{}:\allowbreak{}id2. 2) Ejecutar PATCH /\allowbreak{}api/\allowbreak{}v1/\allowbreak{}orders/\allowbreak{}:\allowbreak{}id1/\allowbreak{}contact. 3) Reconsultar comprobante de P1 y datos de P2.}{P1 refleja contacto corregido; GET /\allowbreak{}orders/\allowbreak{}:\allowbreak{}id1/\allowbreak{}receipt mantiene snapshot original y P2 conserva sus datos y contactRevision sin cambio.}{Comparación antes/\allowbreak{}después demuestra aislamiento total: comprobante congelado y segundo pedido intactos.}
\caso{CP-012}{\textbf{Rechazo de corrección de contacto por rol no autorizado.} NR-03, NR-07; NP-04 (C07.3, Alta)}{Tokens JWT de CLIENT-A y DRIVER-A; orderId válido con contactRevision actual; payload válido de corrección de contacto.}{1) Enviar PATCH /\allowbreak{}api/\allowbreak{}v1/\allowbreak{}orders/\allowbreak{}:\allowbreak{}id/\allowbreak{}contact con token CLIENT-A. 2) Repetir con token DRIVER-A. 3) Consultar contacto y contact-history como ADMIN.}{Ambas peticiones retornan HTTP 403 Forbidden sin ejecutar la corrección.}{Contacto, contactRevision e historial de auditoría permanecen idénticos tras los dos intentos rechazados.}
\caso{CP-013}{\textbf{Horario ADMIN límite inicial 07:\allowbreak{}00 America/\allowbreak{}Guayaquil.} NR-14; NP-06 (C10.5, Alta)}{Token ADMIN; pedido ADMIN válido con scheduledDeliveryDate en fecha futura a las 07:\allowbreak{}00:\allowbreak{}00-05:\allowbreak{}00; Idempotency-Key UUIDv4.}{1) Verificar deliverySchedule 07:\allowbreak{}00-17:\allowbreak{}00 en shop-settings. 2) Enviar POST /\allowbreak{}api/\allowbreak{}v1/\allowbreak{}orders/\allowbreak{}admin con hora 07:\allowbreak{}00-05:\allowbreak{}00. 3) Consultar pedido y comprobante.}{HTTP 201; pedido creado en estado PENDING con hora local 07:\allowbreak{}00 en America/\allowbreak{}Guayaquil y comprobante emitido.}{Límite inicial inclusivo 07:\allowbreak{}00 aceptado y stock descontado correctamente.}
\caso{CP-014}{\textbf{Horario ADMIN límite final 17:\allowbreak{}00 America/\allowbreak{}Guayaquil.} NR-14; NP-06 (C10.6, Alta)}{Token ADMIN; pedido ADMIN válido con scheduledDeliveryDate en fecha futura a las 17:\allowbreak{}00:\allowbreak{}00-05:\allowbreak{}00; Idempotency-Key UUIDv4.}{1) Verificar deliverySchedule 07:\allowbreak{}00-17:\allowbreak{}00. 2) Enviar POST /\allowbreak{}api/\allowbreak{}v1/\allowbreak{}orders/\allowbreak{}admin con hora 17:\allowbreak{}00-05:\allowbreak{}00. 3) Consultar pedido y comprobante.}{HTTP 201; pedido creado en estado PENDING con hora local 17:\allowbreak{}00 en America/\allowbreak{}Guayaquil y comprobante generado.}{Límite final inclusivo 17:\allowbreak{}00 aceptado con reserva de stock y datos consistentes.}
\caso{CP-015}{\textbf{Rechazo horario ADMIN inválido 06:\allowbreak{}59 (422).} NR-14; NP-06 (C10.7, Alta)}{Token ADMIN; pedido ADMIN con scheduledDeliveryDate en fecha futura a las 06:\allowbreak{}59:\allowbreak{}00-05:\allowbreak{}00; Idempotency-Key UUIDv4.}{1) Registrar conteo de pedidos, comprobantes y stock. 2) Enviar POST /\allowbreak{}api/\allowbreak{}v1/\allowbreak{}orders/\allowbreak{}admin con hora 06:\allowbreak{}59-05:\allowbreak{}00. 3) Reconsultar conteos y stock.}{HTTP 422 con error.\allowbreak{}code=\allowbreak{}DELIVERY\_OUTSIDE\_HOURS rechazando el horario previo a 07:\allowbreak{}00.}{Sin pedido ni comprobante creado y stock del producto sin variación.}
\caso{CP-016}{\textbf{Rechazo horario ADMIN inválido 17:\allowbreak{}01 (422).} NR-14; NP-06 (C10.8, Alta)}{Token ADMIN; pedido ADMIN con scheduledDeliveryDate en fecha futura a las 17:\allowbreak{}01:\allowbreak{}00-05:\allowbreak{}00; Idempotency-Key UUIDv4.}{1) Registrar conteo de pedidos, comprobantes y stock. 2) Enviar POST /\allowbreak{}api/\allowbreak{}v1/\allowbreak{}orders/\allowbreak{}admin con hora 17:\allowbreak{}01-05:\allowbreak{}00. 3) Reconsultar conteos y stock.}{HTTP 422 con error.\allowbreak{}code=\allowbreak{}DELIVERY\_OUTSIDE\_HOURS rechazando el horario posterior a 17:\allowbreak{}00.}{Sin efectos parciales: conteos de pedidos, comprobantes y stock idénticos antes y después.}
\caso{CP-017}{\textbf{Agrupación temporal por extremos (10:\allowbreak{}00, 10:\allowbreak{}25, 10:\allowbreak{}50).} NR-18, NR-22; NP-08 (C13.1, Media)}{Token ADMIN; 3 pedidos PENDING elegibles hoy dentro de cobertura con horas 10:\allowbreak{}00, 10:\allowbreak{}25 y 10:\allowbreak{}50 (-05:\allowbreak{}00).}{1) Confirmar routePlanning.\allowbreak{}timeWindow\allowbreak{}Margin\allowbreak{}Minutes=\allowbreak{}30. 2) Invocar GET /\allowbreak{}api/\allowbreak{}v1/\allowbreak{}route-proposals?date=\allowbreak{}HOY como ADMIN. 3) Inspeccionar grupos y paradas.}{HTTP 200 con 2 propuestas separadas: [10:\allowbreak{}00, 10:\allowbreak{}25] y [10:\allowbreak{}50], porque 10:\allowbreak{}50 - 10:\allowbreak{}00 = 50 \ensuremath{>} 30 min (sin encadenamiento).}{Agrupación cumple max - min \ensuremath{\leq} 30 por extremos y orden cronológico de paradas sin persistir rutas.}
\caso{CP-018}{\textbf{Agrupación por extremos con separación de 29 minutos.} NR-18, NR-22; NP-08 (C13.2, Media)}{Token ADMIN; 2 pedidos PENDING elegibles hoy con horas 10:\allowbreak{}00 y 10:\allowbreak{}29 (-05:\allowbreak{}00), diferencia 29 minutos.}{1) Verificar margen=\allowbreak{}30. 2) Ejecutar GET /\allowbreak{}api/\allowbreak{}v1/\allowbreak{}route-proposals?date=\allowbreak{}HOY. 3) Revisar cantidad de rutas y secuencia de paradas.}{HTTP 200 agrupando ambos pedidos en una única propuesta ordenada cronológicamente (10:\allowbreak{}00, luego 10:\allowbreak{}29).}{Separación de 29 min \ensuremath{\leq} 30 se consolida en una ruta sin exceder capacidad máxima de 4.}
\caso{CP-019}{\textbf{Agrupación por extremos con separación exacta de 30 min.} NR-18, NR-22; NP-08 (C13.3, Media)}{Token ADMIN; 2 pedidos PENDING elegibles hoy con horas 10:\allowbreak{}00 y 10:\allowbreak{}30 (-05:\allowbreak{}00), diferencia exacta 30 minutos.}{1) Verificar margen=\allowbreak{}30. 2) Ejecutar GET /\allowbreak{}api/\allowbreak{}v1/\allowbreak{}route-proposals?date=\allowbreak{}HOY. 3) Inspeccionar propuestas y timeWindow\allowbreak{}Margin\allowbreak{}Minutes retornado.}{HTTP 200 agrupando ambos pedidos en una sola ruta al cumplirse la condición inclusiva max - min \ensuremath{\leq} 30.}{Igualdad exacta al límite (30 == 30) incluida en el mismo grupo con orden cronológico.}
\caso{CP-020}{\textbf{Separación de rutas por extremos con 31 minutos.} NR-18, NR-22; NP-08 (C13.4, Media)}{Token ADMIN; 2 pedidos PENDING elegibles hoy con horas 10:\allowbreak{}00 y 10:\allowbreak{}31 (-05:\allowbreak{}00), diferencia 31 minutos.}{1) Verificar margen=\allowbreak{}30. 2) Ejecutar GET /\allowbreak{}api/\allowbreak{}v1/\allowbreak{}route-proposals?date=\allowbreak{}HOY. 3) Verificar separación de grupos en proposals.}{HTTP 200 devolviendo 2 rutas propuestas separadas (1 parada cada una) por exceder el margen de 30 minutos.}{Separación de 31 min (\ensuremath{>} 30) impide la fusión de pedidos en la misma propuesta.}
\caso{CP-021}{\textbf{Cambio de margen horario (5..180) y desempate OSRM.} NR-18, NR-22; NP-08 (C13.5, Media)}{Token ADMIN; PATCH admin/\allowbreak{}shop-settings con routePlanning.\allowbreak{}timeWindow\allowbreak{}Margin\allowbreak{}Minutes=\allowbreak{}60; casos límite inválidos 4 y 181.}{1) Probar márgenes 4 y 181 \ensuremath{\rightarrow} 400. 2) Actualizar margen a 60 vía PATCH /\allowbreak{}api/\allowbreak{}v1/\allowbreak{}admin/\allowbreak{}shop-settings. 3) Consultar GET /\allowbreak{}api/\allowbreak{}v1/\allowbreak{}route-proposals?date=\allowbreak{}HOY.}{4 y 181 rechazados (400); con 60 min los 3 pedidos (10:\allowbreak{}00-10:\allowbreak{}50) se agrupan y aplican desempate OSRM durationsS primario / distancesM secundario.}{Margen validado en rango 5..180, persistido en settings y reflejado en nuevas propuestas con métricas viales coherentes.}
\caso{CP-022}{\textbf{Creación de pedido sin fecha programada (scheduledDeliveryDate nulo).} NR-10; NP-15 (C26.1, Media)}{Token CLIENT-A; payload POST /\allowbreak{}api/\allowbreak{}v1/\allowbreak{}orders con deliveryAddress, items y coordenadas válidas, omitiendo scheduledDeliveryDate.}{1) Enviar POST /\allowbreak{}api/\allowbreak{}v1/\allowbreak{}orders sin scheduledDeliveryDate. 2) Consultar GET /\allowbreak{}api/\allowbreak{}v1/\allowbreak{}orders/\allowbreak{}:\allowbreak{}id.}{HTTP 201; pedido creado con status=\allowbreak{}PENDING y scheduledDeliveryDate=\allowbreak{}null sin asignar fecha automática.}{Contrato opcional preservado: scheduledDeliveryDate permanece null y stock queda reservado.}
\caso{CP-023}{\textbf{Exclusión en propuestas y 409 en asignación sin fecha.} NR-10; NP-15 (C26.2, Media)}{Token ADMIN; pedido PENDING con scheduledDeliveryDate=\allowbreak{}null; fecha operativa HOY; driverId disponible.}{1) Identificar el pedido sin fecha en el listado administrativo. 2) Consultar propuestas de HOY. 3) Intentar confirmar y asignar vía POST /\allowbreak{}routes y POST /\allowbreak{}routes/\allowbreak{}assign.}{El pedido se identifica con fecha nula, no figura en propuestas y ambos intentos de asignación responden 409.}{Pedido identificable, no propuesto ni asignado; sin nueva ruta ni OrderClaim.}
\caso{CP-024}{\textbf{Idempotencia Shop checkout con replay 201.} NR-01, NR-05, NR-10 (C30.1, Alta)}{Token CLIENT-A; quoteToken vigente; Idempotency-Key UUIDv4 fijo y payload idéntico en dos solicitudes consecutivas.}{1) Enviar POST /\allowbreak{}api/\allowbreak{}v1/\allowbreak{}shop/\allowbreak{}checkout con Idempotency-Key K1. 2) Reenviar exactamente el mismo POST con K1 y mismo payload. 3) Consultar pedidos, recibos, stock y saldo.}{Ambas llamadas retornan HTTP 201 con idéntico orderId y receipt.\allowbreak{}id (replay idempotente).}{Solo existe un pedido, un comprobante, un descuento de stock y un débito de saldo demo.}
\caso{CP-025}{\textbf{Recuperación Shop ante timeout con mismo Idempotency-Key.} NR-01, NR-05, NR-10 (C30.2, Alta)}{Token CLIENT-A; quoteToken vigente; Idempotency-Key UUIDv4 K2 con payload válido bajo corte simulado de respuesta HTTP.}{1) Enviar POST /\allowbreak{}api/\allowbreak{}v1/\allowbreak{}shop/\allowbreak{}checkout con K2 provocando corte/\allowbreak{}timeout tras commit en servidor. 2) Reintentar con mismo K2 y payload. 3) Verificar stock y saldo.}{El reintento retorna HTTP 201 con el resultado original sin duplicar la transacción.}{Recuperación consistente: 1 pedido, 1 comprobante, 1 descuento de stock y 1 débito tras corte de red.}
\caso{CP-026}{\textbf{Conflicto Idempotency-Key con payload distinto (409).} NR-01, NR-05, NR-10 (C30.3, Alta)}{Token CLIENT-A; mismo Idempotency-Key UUIDv4 K3 reutilizado con payload alterado (cantidad o dirección distinta).}{1) Registrar conteo de pedidos, recibos, stock y saldo tras K3 original. 2) Enviar POST /\allowbreak{}api/\allowbreak{}v1/\allowbreak{}shop/\allowbreak{}checkout con K3 y payload distinto. 3) Reconsultar estado.}{HTTP 409 con error.\allowbreak{}code=\allowbreak{}IDEMPOTENCY\_KEY\_REUSED rechazando la reutilización conflictiva.}{No se crea segunda venta ni comprobante, ni se altera stock o saldo respecto al estado previo.}
\caso{CP-027}{\textbf{Liberación auditada de DRAFT caducado (RELEASE\_DRAFT).} NR-16, NR-20; NP-07 (C11.5, Alta)}{Token ADMIN; ruta DRAFT de ayer con pedido de ayer; driverId disponible; requestId UUIDv4 y reason válido.}{1) Intentar asignar la ruta y ejecutar RELEASE\_DRAFT. 2) Verificar fecha original, auditoría y replay. 3) Registrar acuerdo mediante CONTACT, aplicar RESCHEDULE con fecha futura y consultar propuestas de esa fecha.}{Asignación responde 409 ROUTE\_DATE\_NOT\_TODAY; se liberan paradas y claim sin cambiar la fecha. Tras acuerdo y reprogramación, el pedido aparece en propuestas de la nueva fecha.}{Pedido PENDING sin claim activo; no se reprograma automáticamente; auditoría y replay idempotente; elegibilidad posterior a la reprogramación.}
\end{disenio}

\clearpage
% FRONTEND WEB - diseño de 27 casos (CP-028..CP-054)
% GUÍA EQUIPO: Responsable asignado: Kerly Chicaiza (Frontend Web).
\paragraph{Módulo Web --- Tienda en línea y panel de rutas (27 casos: CP-028 a CP-054).}

\begin{disenio}
\caso{CP-028}{\textbf{Historial unificado Mis pedidos con compras SHOP y ADMIN propias.} NR-01, NR-07, NP-19 (C05.1, Alta)}{Sesión web de CLIENTE-A; una orden de origen SHOP y una de origen ADMIN del mismo titular.}{1) Abrir Mis pedidos en web. 2) Verificar listado vía GET /\allowbreak{}api/\allowbreak{}v1/\allowbreak{}shop/\allowbreak{}orders. 3) Abrir el detalle de ambas órdenes.}{La lista unificada muestra solo las dos compras propias de CLIENTE-A con su origen (SHOP/\allowbreak{}ADMIN) y estado de seguimiento.}{Coinciden identidades, orígenes y seguimiento entre UI y API; no aparecen órdenes de terceros ni invitadas.}
\caso{CP-029}{\textbf{Bloqueo en lista y detalle ante pedido ajeno o invitado coincidente.} NR-01, NR-07, NP-19 (C05.2, Alta)}{Sesión web de CLIENTE-A; ID de pedido de CLIENTE-B y pedido invitado con teléfono coincidente.}{1) Revisar Mis pedidos de CLIENTE-A. 2) Intentar abrir en web y GET /\allowbreak{}api/\allowbreak{}v1/\allowbreak{}shop/\allowbreak{}orders/\allowbreak{}:\allowbreak{}id el pedido ajeno y el invitado.}{El pedido invitado no se vincula por teléfono en la lista; ambos accesos por ID responden 404 ORDER\_NOT\_FOUND sin datos.}{Sin fuga de información en UI ni API; el código 404 ORDER\_NOT\_FOUND bloquea detalle ajeno e invitado.}
\caso{CP-030}{\textbf{Pedido sin PurchaseReceipt y sin importes inventados.} NR-05, NR-08, NP-03 (C08.1, Media)}{Sesión web de CLIENTE-A; ID de pedido previo propio con receiptAvailable=\allowbreak{}false.}{1) Consultar lista y detalle del pedido en web y GET /\allowbreak{}api/\allowbreak{}v1/\allowbreak{}shop/\allowbreak{}orders. 2) Solicitar /\allowbreak{}api/\allowbreak{}v1/\allowbreak{}shop/\allowbreak{}orders/\allowbreak{}:\allowbreak{}id/\allowbreak{}receipt.}{Lista y detalle marcan receiptAvailable=\allowbreak{}false e informan comprobante no disponible; /\allowbreak{}receipt responde 404 RECEIPT\_UNAVAILABLE.}{No se reconstruye comprobante ni se inventan importes o impuestos en UI ni en API; datos almacenados intactos.}
\caso{CP-031}{\textbf{Cobertura tienda: punto interior al polígono aceptado en web.} NR-11, NR-12, NR-13, NP-05 (C09.1, Alta)}{Sesión CLIENTE-A; coordenadas sintéticas dentro del polígono de cobertura vigente; carrito Shop válido.}{1) Seleccionar punto interior en el mapa de checkout web. 2) Verificar POST /\allowbreak{}api/\allowbreak{}v1/\allowbreak{}store-coverage/\allowbreak{}evaluate. 3) Completar compra.}{La evaluación retorna position=\allowbreak{}INTERIOR y accepted=\allowbreak{}true; la interfaz permite avanzar y confirmar el pedido Shop.}{El punto interior queda aceptado sin bloqueos y la orden se registra con sus coordenadas válidas.}
\caso{CP-032}{\textbf{Cobertura tienda: punto exactamente en el borde del polígono.} NR-11, NR-12, NR-13, NP-05 (C09.2, Alta)}{Sesión CLIENTE-A; coordenadas exactamente sobre el borde del polígono de StoreCoverage.}{1) Ingresar punto en el borde en checkout web. 2) Evaluar vía POST /\allowbreak{}api/\allowbreak{}v1/\allowbreak{}store-coverage/\allowbreak{}evaluate. 3) Confirmar si permite compra.}{La evaluación devuelve position=\allowbreak{}BOUNDARY y accepted=\allowbreak{}true, habilitando la compra en la web.}{El punto en el borde es aceptado consistentemente en UI y API según la regla geométrica vigente.}
\caso{CP-033}{\textbf{Cobertura tienda: punto exterior rechazado con 422 en web.} NR-11, NR-12, NR-13, NP-05 (C09.3, Alta)}{Sesión CLIENTE-A; coordenadas fuera del polígono de cobertura; producto disponible en carrito.}{1) Seleccionar punto exterior en checkout web. 2) Consultar evaluate e intentar confirmar compra. 3) Verificar stock y órdenes.}{Evaluate retorna EXTERIOR y accepted=\allowbreak{}false; el intento de compra se rechaza con 422 DELIVERY\_OUTSIDE\_COVERAGE y aviso en UI.}{No se crea pedido ni comprobante, y el stock y saldo simulado del cliente permanecen sin cambios.}
\caso{CP-034}{\textbf{Cambio de cobertura A a B conserva pedidos ya aceptados.} NR-11, NR-12, NR-13, NP-05 (C09.4, Alta)}{Polígono A inicial con pedido aceptado en punto P; polígono B nuevo que excluye al punto P.}{1) Como ADMIN actualizar polígono a B en panel web (PUT /\allowbreak{}api/\allowbreak{}v1/\allowbreak{}admin/\allowbreak{}store-coverage). 2) Consultar el pedido previo.}{El cambio de cobertura se guarda, pero el pedido aceptado en P conserva su estado, dirección y datos sin cancelación.}{Ningún pedido previo aceptado es cancelado, alterado ni eliminado automáticamente tras el cambio de polígono.}
\caso{CP-035}{\textbf{Compra nueva tras cambio de cobertura aplica área B vigente.} NR-11, NR-12, NR-13, NP-05 (C09.5, Alta)}{Sesión CLIENTE-A; punto P válido en área A pero exterior al área B vigente; carrito válido.}{1) Intentar nueva compra web para punto P bajo área B. 2) Verificar respuesta y estado de carrito/\allowbreak{}pedidos.}{La compra nueva se rechaza en UI y API con 422 DELIVERY\_OUTSIDE\_COVERAGE al aplicar el polígono B vigente.}{Sin pedido, comprobante ni cobro nuevos; el pedido ya aceptado conserva su dirección original.}
\caso{CP-036}{\textbf{Horario Shop límite inicial 07:\allowbreak{}00 America/\allowbreak{}Guayaquil aceptado.} NR-14, NP-06 (C10.1, Alta)}{Sesión CLIENTE-A; fecha futura válida con hora 07:\allowbreak{}00 (-05:\allowbreak{}00 America/\allowbreak{}Guayaquil); pago DEMO\_APPROVED.}{1) Seleccionar 07:\allowbreak{}00 en checkout web. 2) Confirmar compra vía POST /\allowbreak{}api/\allowbreak{}v1/\allowbreak{}shop/\allowbreak{}checkout. 3) Consultar pedido.}{El horario 07:\allowbreak{}00 se acepta en UI y backend; la orden se crea con la hora programada correspondiente.}{HTTP 201 en checkout, pedido y comprobante generados y hora almacenada coherente con 07:\allowbreak{}00 America/\allowbreak{}Guayaquil.}
\caso{CP-037}{\textbf{Horario Shop límite final 17:\allowbreak{}00 America/\allowbreak{}Guayaquil aceptado.} NR-14, NP-06 (C10.2, Alta)}{Sesión CLIENTE-A; fecha futura válida con hora 17:\allowbreak{}00 (-05:\allowbreak{}00 America/\allowbreak{}Guayaquil); carrito y cobertura válidos.}{1) Seleccionar 17:\allowbreak{}00 en checkout web. 2) Confirmar la compra y consultar el detalle del pedido.}{El horario límite 17:\allowbreak{}00 es aceptado en web y API, creando la orden con entrega programada a las 17:\allowbreak{}00.}{Compra confirmada (201) sin rechazos indebidos en el extremo superior 17:\allowbreak{}00.}
\caso{CP-038}{\textbf{Horario Shop inválido 06:\allowbreak{}59 rechazado sin orden ni cobro.} NR-14, NP-06 (C10.3, Alta)}{Sesión CLIENTE-A; fecha futura válida con hora 06:\allowbreak{}59 (-05:\allowbreak{}00 America/\allowbreak{}Guayaquil); tarjeta DEMO\_APPROVED.}{1) Intentar confirmar compra para las 06:\allowbreak{}59 en web o solicitud checkout. 2) Verificar stock, saldo y órdenes.}{Se rechaza con 422 DELIVERY\_OUTSIDE\_HOURS y mensaje claro en UI; no se genera orden ni comprobante.}{Stock y saldo simulado intactos tras el rechazo; cero órdenes creadas para 06:\allowbreak{}59.}
\caso{CP-039}{\textbf{Horario Shop inválido 17:\allowbreak{}01 rechazado sin orden ni cobro.} NR-14, NP-06 (C10.4, Alta)}{Sesión CLIENTE-A; fecha futura válida con hora 17:\allowbreak{}01 (-05:\allowbreak{}00 America/\allowbreak{}Guayaquil); carrito válido.}{1) Intentar confirmar compra para las 17:\allowbreak{}01 en checkout web. 2) Verificar órdenes, stock y saldo.}{El servidor responde 422 DELIVERY\_OUTSIDE\_HOURS y la web muestra error sin crear pedido ni realizar cobro.}{Sin orden ni comprobante emitidos; inventario y saldo del cliente permanecen sin variación.}
\caso{CP-040}{\textbf{Cambio de horario en shop-settings genera 409 al confirmar.} NR-14, NP-06 (C10.9, Alta)}{Cotización Shop abierta con horario 07:\allowbreak{}00-17:\allowbreak{}00; ADMIN cambia ventana vía PATCH /\allowbreak{}api/\allowbreak{}v1/\allowbreak{}admin/\allowbreak{}shop-settings.}{1) Cotizar en web como CLIENTE-A. 2) Como ADMIN actualizar horario. 3) Confirmar compra pendiente en web.}{El checkout responde 409 SHOP\_SETTINGS\_CHANGED; la web informa el cambio y exige recotizar sin confirmar con reglas previas.}{No se crea orden ni cobro bajo el horario obsoleto; tras recotizar aplican las condiciones vigentes.}
\caso{CP-041}{\textbf{Verificación UI web de aviso de hora de entrega aproximada.} NR-14, NP-06 (C10.10, Alta)}{Sesión CLIENTE-A; pantalla de selección de horario en checkout Shop y detalle de pedido creado.}{1) Inspeccionar textos de horario en checkout web. 2) Completar compra y revisar textos en detalle del pedido.}{La interfaz informa claramente que la hora de entrega es aproximada y no promete exactitud ni margen de +-30 min.}{El aviso visible es explícito sobre carácter aproximado y no garantiza ventana +-30 min (validación puramente UI).}
\caso{CP-042}{\textbf{Propuesta de rutas en panel ADMIN es solo lectura sin persistir.} NR-20, NP-07, NP-13 (C14.1, Alta)}{Sesión ADMIN en panel web; pedidos elegibles de hoy (isToday=\allowbreak{}true) y otro día (isToday=\allowbreak{}false).}{1) Solicitar propuesta en panel (GET /\allowbreak{}api/\allowbreak{}v1/\allowbreak{}route-proposals) para hoy y fecha futura. 2) Consultar rutas persistidas.}{Devuelve propuestas temporales con isToday/\allowbreak{}canConfirm correctos; consultar rutas persistidas muestra cero rutas creadas.}{Ninguna ruta ni reserva se persiste en BD al generar propuestas; canConfirm=\allowbreak{}false en fechas distintas de hoy.}
\caso{CP-043}{\textbf{Confirmación de borrador DRAFT en hoy persiste ruta sin conductor.} NR-20, NP-07, NP-13 (C14.2, Alta)}{Sesión ADMIN; propuesta del día actual con canConfirm=\allowbreak{}true y pedidos elegibles.}{1) Confirmar propuesta en panel ADMIN (POST /\allowbreak{}api/\allowbreak{}v1/\allowbreak{}routes). 2) Consultar listado de rutas y reservas de pedidos.}{Se persiste una ruta en estado DRAFT sin conductor asignado (driverId=\allowbreak{}null) y visible en el panel.}{Estado DRAFT confirmado sin duplicar rutas ni asignar conductor automáticamente; pedidos reclamados por la ruta.}
\caso{CP-044}{\textbf{Asignación separada de ruta DRAFT a conductor disponible.} NR-20, NP-07, NP-13 (C14.3, Alta)}{Sesión ADMIN; ID de ruta DRAFT y REPARTIDOR-A sin ausencias ni reservas activas.}{1) Desde panel ADMIN asignar REPARTIDOR-A (POST /\allowbreak{}api/\allowbreak{}v1/\allowbreak{}routes/\allowbreak{}:\allowbreak{}id/\allowbreak{}assign). 2) Consultar ruta y disponibilidad.}{La ruta pasa a estado ASSIGNED con REPARTIDOR-A y crea reserva activa del conductor en el sistema.}{Transición DRAFT\ensuremath{\rightarrow}ASSIGNED persistida como acción separada; reserva creada y reflejada en disponibilidad.}
\caso{CP-045}{\textbf{Rechazo de asignación a conductor no disponible o ausente (409).} NR-20, NP-07, NP-13 (C14.4, Alta)}{Sesión ADMIN; ruta DRAFT sin conductor y REPARTIDOR-B con ausencia vigente o reserva previa.}{1) Intentar asignar REPARTIDOR-B a la ruta DRAFT en panel ADMIN. 2) Consultar rutas sin conductor y reservas.}{La asignación se rechaza con HTTP 409 (DRIVER\_ABSENT o conflicto de reserva) y aviso claro en UI.}{Sin asignación parcial ni reserva indebida; la ruta permanece en DRAFT sin conductor y visible en el panel.}
\caso{CP-046}{\textbf{Motor vial OSRM indisponible muestra aviso sin rutas en línea recta.} NR-19, NR-21, NP-09 (C16.1, Alta)}{Sesión ADMIN; pedidos elegibles de hoy con motor vial OSRM detenido o inaccesible.}{1) Solicitar propuestas en panel ADMIN (GET /\allowbreak{}api/\allowbreak{}v1/\allowbreak{}route-proposals). 2) Observar aviso UI y respuesta HTTP.}{El servidor responde 503 ROAD\_ENGINE\_UNAVAILABLE (o 502) y el panel muestra aviso claro sin generar rutas en línea recta.}{Cero propuestas ficticias o trazados en línea recta presentados como resultado vial; nada se persiste.}
\caso{CP-047}{\textbf{Preservación de pedidos pendientes y rutas tras falla OSRM.} NR-19, NR-21, NP-09 (C16.2, Alta)}{Sesión ADMIN; estado previo de pedidos PENDING y ruta existente antes y después del error 503/\allowbreak{}502.}{1) Consultar listado de pedidos pendientes y rutas existentes en panel tras el rechazo 503/\allowbreak{}502.}{Los pedidos siguen en PENDING y las rutas previamente confirmadas o asignadas se conservan intactas.}{Comparación antes/\allowbreak{}después confirma ausencia de mutaciones, cancelaciones o reservas parciales tras el error vial.}
\caso{CP-048}{\textbf{Recuperación tras indisponibilidad OSRM genera propuestas válidas.} NR-19, NR-21, NP-09 (C16.3, Alta)}{Sesión ADMIN; servicio OSRM nuevamente operativo y pedidos pendientes del día.}{1) Verificar estado en panel/\allowbreak{}road-engine. 2) Solicitar nuevamente propuestas desde panel ADMIN.}{El panel genera propuestas viales válidas con rutas, distancias y duraciones estimadas normalmente.}{HTTP 200 en propuestas tras recuperación; resultados coherentes con el motor vial sin errores residuales.}
\caso{CP-049}{\textbf{Origen y retorno común en tienda única reflejado en propuesta y mapa.} NR-11, NR-19, NP-08 (C17.1, Media)}{Sesión ADMIN; pedidos del día con destino válido dentro de cobertura y tienda única configurada.}{1) Solicitar propuesta con geometría en panel ADMIN. 2) Verificar punto de inicio y fin en JSON y mapa.}{La propuesta y el mapa web parten y regresan a la ubicación única de StoreCoverage.\allowbreak{}location, sin depósitos por zona.}{Coordenadas de origen y retorno coinciden exactamente con StoreCoverage.\allowbreak{}location configurada.}
\caso{CP-050}{\textbf{Paradas y trazado vial en panel ADMIN coinciden con motor vial.} NR-11, NR-19, NP-08 (C17.2, Media)}{Sesión ADMIN; propuesta con 2 a 4 paradas en horarios compatibles y trazado geométrico OSRM.}{1) Abrir mapa de la propuesta en panel ADMIN. 2) Comparar secuencia de paradas y polilínea con respuesta API.}{El mapa muestra las paradas en el orden cronológico devuelto y sigue el trazado vial OSRM, no rectas entre puntos.}{Orden de paradas y geometría renderizada corresponden fielmente con los datos viales recibidos del backend.}
\caso{CP-051}{\textbf{Distancias y duraciones mostradas como estimadas sin tráfico en vivo.} NR-11, NR-19, NP-08 (C17.3, Media)}{Sesión ADMIN; metadatos roadMetrics (estimate=\allowbreak{}true, trafficAware=\allowbreak{}false) y resumen mostrado en UI.}{1) Revisar métricas e indicadores en panel de propuestas. 2) Contrastar textos visibles con roadMetrics del API.}{La UI presenta distancias y duraciones como estimadas (estimate=\allowbreak{}true, trafficAware=\allowbreak{}false) sin afirmar ETA ni tráfico vivo.}{Ninguna etiqueta web promete hora exacta de llegada, tráfico en tiempo real ni solución óptima global.}
\caso{CP-052}{\textbf{Regresión Shop: compra válida con DEMO\_APPROVED crea pedido y recibo.} NR-01, NR-05, NR-10 (C27.1, Alta)}{Sesión CLIENTE-A; producto disponible, cobertura y horario válidos, tarjeta simulada DEMO\_APPROVED.}{1) Añadir producto desde catálogo Shop web. 2) Completar checkout con DEMO\_APPROVED. 3) Abrir pedido y comprobante.}{Se crea el pedido y su comprobante simulado; el stock y el saldo demo se descuentan exactamente una vez.}{Comprobante coincide con precios e IVA del servidor; sin exponer número completo ni CVV de tarjeta en evidencias.}
\caso{CP-053}{\textbf{Regresión Shop: carrito mixto con agotado rechazado sin efectos.} NR-01, NR-05, NR-10 (C27.2, Alta)}{Sesión CLIENTE-A; carrito con 1 producto disponible y 1 agotado; stock y saldo inicial medidos.}{1) Intentar cotizar o confirmar carrito mixto en Shop web/\allowbreak{}API. 2) Verificar órdenes, stock disponible y saldo.}{El catálogo marca el agotado como no disponible y el checkout rechaza el carrito mixto sin crear orden parcial.}{Stock del producto disponible y saldo demo quedan idénticos al estado previo; ningún pedido o recibo creado.}
\caso{CP-054}{\textbf{Regresión Shop: precio o total manipulado en cliente no es aceptado.} NR-01, NR-05, NR-10 (C27.3, Alta)}{Sesión CLIENTE-A; payload hacia POST /\allowbreak{}api/\allowbreak{}v1/\allowbreak{}shop/\allowbreak{}checkout con campos price/\allowbreak{}total alterados a valor menor.}{1) Enviar solicitud manipulada interceptada desde cliente. 2) Inspeccionar respuesta HTTP, orden y cobro.}{El servidor rechaza campos no permitidos (400/\allowbreak{}409) o aplica únicamente sus precios reales; nunca cobra el monto alterado.}{Si rechaza no deja orden ni débito; si acepta solo cobra e imprime el importe calculado por el servidor.}
\end{disenio}

\clearpage
% MÓVIL - diseño de 27 casos (CP-055..CP-081)
% GUÍA EQUIPO: Responsable asignado: Carlos Ramos (Móvil).
\paragraph{Módulo Móvil --- Seguridad y operación de entregas (27 casos: CP-055 a CP-081).}

\begin{disenio}
\caso{CP-055}{\textbf{Acceso sin sesión JWT a rutas, pedidos y configuración.} NR-01, NR-03, NR-07, NR-23; NP-14 (C06.1, Alta)}{Solicitudes HTTP sin cabecera Authorization a GET /\allowbreak{}api/\allowbreak{}v1/\allowbreak{}drivers/\allowbreak{}me/\allowbreak{}routes, GET /\allowbreak{}api/\allowbreak{}v1/\allowbreak{}shop/\allowbreak{}orders y GET /\allowbreak{}api/\allowbreak{}v1/\allowbreak{}admin/\allowbreak{}shop-settings.}{1) Invocar los tres endpoints sin token JWT. 2) Inspeccionar códigos HTTP y cuerpos de respuesta.}{Todas las peticiones se rechazan con HTTP 401 sin devolver rutas, pedidos ni parámetros de configuración.}{100 \% de los endpoints protegidos responden 401 y sus respuestas no filtran datos internos ni listados.}
\caso{CP-056}{\textbf{Aislamiento entre CLIENT-B y pedido/\allowbreak{}comprobante de CLIENT-A.} NR-01, NR-03, NR-07, NR-23; NP-14 (C06.2, Alta)}{JWT de CLIENT-B y UUID de pedido/\allowbreak{}comprobante perteneciente a CLIENT-A.}{1) Consultar GET /\allowbreak{}api/\allowbreak{}v1/\allowbreak{}shop/\allowbreak{}orders/\allowbreak{}:\allowbreak{}id y /\allowbreak{}receipt como CLIENT-B. 2) Intentar mutar o consultar /\allowbreak{}orders/\allowbreak{}:\allowbreak{}id ajeno.}{El backend rechaza cada intento con HTTP 403 o 404 sin exponer montos, destinatario ni comprobante de CLIENT-A.}{Ninguna respuesta revela datos de CLIENT-A y el pedido original permanece sin cambios.}
\caso{CP-057}{\textbf{Aislamiento de DRIVER-B ante ruta y parada de DRIVER-A.} NR-01, NR-03, NR-07, NR-23; NP-14 (C06.3, Alta)}{Sesión JWT de DRIVER-B y routeId/\allowbreak{}stopId pertenecientes a DRIVER-A.}{1) Abrir Mi Ruta en móvil con DRIVER-B. 2) Consultar ruta ajena e intentar iniciar, fallar o cerrar parada/\allowbreak{}ruta de DRIVER-A.}{Móvil no lista rutas ajenas y la API rechaza consultas y mutaciones sobre DRIVER-A con HTTP 403 o 404.}{Ruta y paradas de DRIVER-A conservan su estado intacto sin fugas de información a DRIVER-B.}
\caso{CP-058}{\textbf{Registro público con rol ADMIN rechazado sin escalación.} NR-01, NR-03, NR-07, NR-23; NP-14 (C06.4, Alta)}{Payload de registro con name, email nuevo válido, password válido y campo adicional role: ADMIN.}{1) Enviar POST /\allowbreak{}api/\allowbreak{}v1/\allowbreak{}auth/\allowbreak{}register con role ADMIN. 2) Intentar iniciar sesión con esas credenciales y verificar usuarios.}{El servidor rechaza la petición con HTTP 400 o 403 y no crea cuenta ADMIN ni usuario privilegiado.}{El intento de login posterior falla y no existe ninguna cuenta registrada con el correo enviado.}
\caso{CP-059}{\textbf{Intento de CLIENT de modificar configuración o cobertura.} NR-01, NR-03, NR-07, NR-23; NP-14 (C06.5, Alta)}{JWT de CLIENT-A con payload válido para PATCH /\allowbreak{}api/\allowbreak{}v1/\allowbreak{}admin/\allowbreak{}shop-settings y cambio en /\allowbreak{}api/\allowbreak{}v1/\allowbreak{}admin/\allowbreak{}store-coverage.}{1) Leer configuración actual como ADMIN. 2) Enviar mutaciones de settings y cobertura como CLIENT-A. 3) Releer configuración.}{Ambas peticiones de CLIENT-A se rechazan con HTTP 403 y los valores de tienda y cobertura no cambian.}{Comparación antes/\allowbreak{}después como ADMIN confirma configuración y polígono exactamente iguales.}
\caso{CP-060}{\textbf{Acceso autorizado por titular, ADMIN y DRIVER sobre recursos propios.} NR-01, NR-03, NR-07, NR-23; NP-14 (C06.6, Alta)}{JWT válidos de CLIENT-A (pedido propio), ADMIN (configuración/\allowbreak{}despacho) y DRIVER-A (ruta propia).}{1) CLIENT-A lee su pedido/\allowbreak{}comprobante. 2) ADMIN consulta settings y despacho. 3) DRIVER-A consulta sus rutas en móvil/\allowbreak{}API.}{Cada actor accede con HTTP 200 únicamente a sus recursos permitidos sin exposición de datos de terceros.}{Respuestas corresponden exactamente a los recursos autorizados de cada rol según contrato.}
\caso{CP-061}{\textbf{Pedido de hoy elegible en propuesta, ruta y asignación.} NR-16, NR-20; NP-07, NP-13 (C11.1, Alta)}{Pedido con scheduledDeliveryDate del día operativo actual en America/\allowbreak{}Guayaquil y DRIVER-A disponible.}{1) Consultar GET /\allowbreak{}api/\allowbreak{}v1/\allowbreak{}route-proposals. 2) Crear ruta con POST /\allowbreak{}api/\allowbreak{}v1/\allowbreak{}routes. 3) Asignar DRIVER-A con POST /\allowbreak{}routes/\allowbreak{}:\allowbreak{}id/\allowbreak{}assign.}{La propuesta incluye el pedido con isToday=\allowbreak{}true, permite crear borrador y asignar la ruta a DRIVER-A.}{Propuesta, ruta DRAFT/\allowbreak{}ASSIGNED y asignación quedan persistidas para la fecha operativa actual.}
\caso{CP-062}{\textbf{Pedido futuro admitido en compra y bloqueado en ruta/\allowbreak{}assign.} NR-16, NR-20; NP-07, NP-13 (C11.2, Alta)}{Pedido válido con fecha programada de mañana en POST /\allowbreak{}api/\allowbreak{}v1/\allowbreak{}shop/\allowbreak{}checkout y POST /\allowbreak{}api/\allowbreak{}v1/\allowbreak{}orders/\allowbreak{}admin.}{1) Crear pedidos futuros. 2) Consultar GET /\allowbreak{}api/\allowbreak{}v1/\allowbreak{}route-proposals?date=\allowbreak{}mañana. 3) Intentar POST /\allowbreak{}routes y /\allowbreak{}assign para mañana.}{Las compras futuras se crean (201) y permiten preview, pero POST /\allowbreak{}routes y assign se rechazan con 409 ROUTE\_DATE\_NOT\_TODAY.}{No se crean rutas ni asignaciones futuras; pedidos permanecen pendientes para su fecha programada.}
\caso{CP-063}{\textbf{Ausencia registrada hoy impide asignar repartidor.} NR-16, NR-20; NP-07, NP-13 (C11.3, Alta)}{Registro en POST /\allowbreak{}api/\allowbreak{}v1/\allowbreak{}drivers/\allowbreak{}:\allowbreak{}driverId/\allowbreak{}absences con startDate/\allowbreak{}endDate de hoy y ruta DRAFT lista.}{1) Registrar ausencia de hoy para DRIVER-B. 2) Consultar disponibilidad. 3) Intentar asignar DRIVER-B en POST /\allowbreak{}routes/\allowbreak{}:\allowbreak{}id/\allowbreak{}assign.}{Disponibilidad reporta isAvailable=\allowbreak{}false y el intento de asignación se rechaza sin asignar la ruta.}{La ruta permanece sin conductor (DRAFT) y no se genera reserva ni salida para DRIVER-B.}
\caso{CP-064}{\textbf{Ausencia futura no altera ni reasigna rutas activas de hoy.} NR-16, NR-20; NP-07, NP-13 (C11.4, Alta)}{Ausencia en POST /\allowbreak{}api/\allowbreak{}v1/\allowbreak{}drivers/\allowbreak{}:\allowbreak{}driverId/\allowbreak{}absences con fechas futuras (mañana o posterior) para DRIVER-A.}{1) Consultar ruta activa de hoy de DRIVER-A. 2) Registrar ausencia futura para DRIVER-A. 3) Reconsultar ruta y paradas de hoy.}{La ausencia futura queda registrada, y la ruta asignada o iniciada de hoy mantiene mismo conductor y paradas.}{No ocurre reasignación automática, cancelación ni alteración de reservas en la ruta de hoy.}
\caso{CP-065}{\textbf{Pedido urgente con ruta IN\_PROGRESS no altera ruta iniciada.} NR-20; NP-09 (C15.1, Alta)}{Pedido nuevo del mismo día registrado tras iniciar R1; consulta a GET /\allowbreak{}api/\allowbreak{}v1/\allowbreak{}dispatch/\allowbreak{}orders/\allowbreak{}:\allowbreak{}id/\allowbreak{}proposal.}{1) Registrar snapshot de R1. 2) Crear pedido urgente para hoy. 3) Consultar propuesta de urgencia e inbox; reconsultar R1.}{La propuesta sugiere NEW\_ROUTE o BLOCKED sin proponer inserción en R1; R1 permanece intacta.}{Paradas, orden y reserva de R1 son idénticos antes y después; no hay inserción automática.}
\caso{CP-066}{\textbf{Decisión manual ADMIN en urgencias: INSERT con cupo o NEW\_ROUTE.} NR-20; NP-09 (C15.2, Alta)}{Pedido urgente de hoy, respuesta de GET /\allowbreak{}dispatch/\allowbreak{}orders/\allowbreak{}:\allowbreak{}id/\allowbreak{}proposal y confirmación explícita en POST /\allowbreak{}dispatch/\allowbreak{}interventions.}{1) Consultar propuesta de urgencia. 2) Confirmar manualmente INSERT (cupo \ensuremath{<}4) o NEW\_ROUTE. 3) Verificar rutas e historial auditado.}{Solo tras confirmación explícita ADMIN se inserta en la ruta elegible (\ensuremath{<}4 paradas) o se crea nueva ruta con auditoría.}{Ninguna ruta supera 4 paradas, el cambio queda auditado y rutas no seleccionadas permanecen intactas.}
\caso{CP-067}{\textbf{Entrega fallida en móvil conserva parada e historial pendiente.} NR-23, NR-24; NP-10 (C18.1, Alta)}{Token DRIVER-A, stopId pendiente y motivo RECIPIENT\_ABSENT o RECIPIENT\_REFUSED en POST /\allowbreak{}api/\allowbreak{}v1/\allowbreak{}routes/\allowbreak{}:\allowbreak{}id/\allowbreak{}stops/\allowbreak{}:\allowbreak{}stopId/\allowbreak{}failure.}{1) Registrar entrega fallida desde móvil con motivo válido. 2) Consultar detalle de ruta, parada y seguimiento.}{Parada pasa a FAILED conservando intento/\allowbreak{}motivo en historial y el pedido queda en estado pendiente de gestión.}{No marca DELIVERED ni CANCELLED, conserva la parada original y el intento auditado.}
\caso{CP-068}{\textbf{Registro de fallo por actor o conductor no autorizado rechazado.} NR-23, NR-24; NP-10 (C18.2, Alta)}{Tokens de DRIVER-B y CLIENT-A sobre POST /\allowbreak{}api/\allowbreak{}v1/\allowbreak{}routes/\allowbreak{}:\allowbreak{}id/\allowbreak{}stops/\allowbreak{}:\allowbreak{}stopId/\allowbreak{}failure de DRIVER-A con motivo válido.}{1) Enviar solicitud de fallo como DRIVER-B y como CLIENT-A. 2) Consultar parada y pedido después de cada intento.}{Cada intento no autorizado se rechaza con HTTP 403 o 409 sin registrar intento fallido.}{La parada sigue PENDING y el pedido no muestra transición ni historial modificado.}
\caso{CP-069}{\textbf{Gestión post-fallo ADMIN: CONTACT y RESCHEDULE con nueva fecha.} NR-23, NR-24; NP-10 (C18.3, Alta)}{Acciones CONTACT y RESCHEDULE en POST /\allowbreak{}api/\allowbreak{}v1/\allowbreak{}dispatch/\allowbreak{}interventions con nueva fecha/\allowbreak{}hora acordada válida.}{1) Registrar CONTACT tras contacto al cliente. 2) Ejecutar RESCHEDULE con fecha/\allowbreak{}hora nueva. 3) Reintentar o consultar historial.}{El pedido queda reprogramado con la nueva fecha/\allowbreak{}hora y conserva el historial de fallo previo en la ruta original.}{Sin duplicar asignación activa ni borrar el intento previo; nueva fecha queda lista para planificación.}
\caso{CP-070}{\textbf{Gestión post-fallo ADMIN: acuerdo de retiro en tienda (PICKUP\_AGREE).} NR-23, NR-24; NP-10 (C18.4, Alta)}{Payload con action PICKUP\_AGREE en POST /\allowbreak{}api/\allowbreak{}v1/\allowbreak{}dispatch/\allowbreak{}interventions para el pedido fallido.}{1) Enviar PICKUP\_AGREE como ADMIN. 2) Consultar seguimiento del pedido como cliente y en despacho.}{El pedido pasa a estado pendiente de retiro en tienda sin marcarlo aún como entregado ni reponer stock.}{Seguimiento muestra pendiente de retiro y conserva registro histórico de la visita fallida.}
\caso{CP-071}{\textbf{Confirmación de retiro (PICKUP\_CONFIRM) y rechazo de repetición.} NR-23, NR-24; NP-10 (C18.5, Alta)}{Acción PICKUP\_CONFIRM en POST /\allowbreak{}api/\allowbreak{}v1/\allowbreak{}dispatch/\allowbreak{}interventions y repetición inmediata de PICKUP\_CONFIRM.}{1) Confirmar retiro físico con PICKUP\_CONFIRM. 2) Repetir PICKUP\_CONFIRM sobre el mismo pedido. 3) Consultar estado e historial.}{La primera confirmación marca entrega por retiro; la repetición se rechaza con 409 PICKUP\_ALREADY\_CONFIRMED.}{Existe una sola entrega por retiro sin registros duplicados ni movimientos adicionales de inventario.}
\caso{CP-072}{\textbf{Recorrido E2E: creación de compra y verificación de estado inicial.} NR-01, NR-05, NR-07, NR-19--NR-21, NR-23; Transversal (C25.1, Alta)}{Carrito válido de CLIENT-A con fecha de hoy, dirección cubierta y medio de pago demo aprobado.}{1) Completar checkout en Shop. 2) Abrir historial/\allowbreak{}seguimiento del cliente y consultar pedido y comprobante.}{Se crea pedido en estado inicial pendiente y comprobante congelado visibles en el seguimiento de CLIENT-A.}{Importes, stock descontado, comprobante y estado inicial coinciden exactamente con el backend.}
\caso{CP-073}{\textbf{Recorrido E2E: propuesta, confirmación y asignación en día operativo.} NR-01, NR-05, NR-07, NR-19--NR-21, NR-23; Transversal (C25.2, Alta)}{Sesión ADMIN, pedido de hoy elegible y DRIVER-A disponible sin ausencias ni rutas activas.}{1) Generar propuesta en web/\allowbreak{}API. 2) Confirmar borrador (DRAFT). 3) Asignar ruta a DRIVER-A y verificar en app móvil.}{Borrador se persiste, queda asignado (ASSIGNED) a DRIVER-A y aparece en Mi Ruta de la app móvil.}{Pedido, parada, reserva del conductor y lista en móvil están sincronizados sin duplicados.}
\caso{CP-074}{\textbf{Recorrido E2E: progreso móvil reflejado en seguimiento del cliente.} NR-01, NR-05, NR-07, NR-19--NR-21, NR-23; Transversal (C25.3, Alta)}{Sesión DRIVER-A en app móvil y sesión CLIENT-A en seguimiento web del mismo pedido.}{1) DRIVER-A inicia ruta en móvil. 2) Marca parada en camino y luego entregada. 3) CLIENT-A refresca seguimiento.}{Cada transición persistida (IN\_PROGRESS, IN\_TRANSIT, DELIVERED) se refleja en tiempo real o consulta en seguimiento cliente.}{Estados visibles en cliente y móvil coinciden con los persistidos en backend sin ubicaciones GPS inventadas.}
\caso{CP-075}{\textbf{Recorrido E2E: desconexión móvil y recuperación al reconectar.} NR-01, NR-05, NR-07, NR-19--NR-21, NR-23; Transversal (C25.4, Alta)}{Desconexión de red controlada en móvil antes/\allowbreak{}durante consulta o acción de parada y reconexión posterior.}{1) Cortar red en móvil y observar aviso ante intento de acción. 2) Reconectar red. 3) Reintentar o refrescar y verificar API.}{Sin red no finge éxito; al reconectar recupera el estado real persistido sin duplicar entregas ni inventar cola offline.}{Backend registra una sola transición válida y el móvil refleja exactamente el estado del servidor.}
\caso{CP-076}{\textbf{Cierre de ruta móvil rechazado con visita aún PENDING.} NR-11, NR-20, NR-23; NP-13, NP-14 (C28.1, Alta)}{Token DRIVER-A y ruta con parada pendiente hacia POST /\allowbreak{}api/\allowbreak{}v1/\allowbreak{}routes/\allowbreak{}:\allowbreak{}id/\allowbreak{}complete con cuerpo \{\}.}{1) Verificar que existe parada PENDING. 2) Intentar cerrar ruta desde móvil o POST /\allowbreak{}routes/\allowbreak{}:\allowbreak{}id/\allowbreak{}complete.}{El servidor rechaza el cierre con HTTP 409 y la ruta continúa en estado IN\_PROGRESS.}{No se registra completedAt, la ruta no pasa a COMPLETED y el conductor sigue ocupado.}
\caso{CP-077}{\textbf{Cierre de ruta COMPLETED con visitas resueltas y una FAILED.} NR-11, NR-20, NR-23; NP-13, NP-14 (C28.2, Alta)}{Token DRIVER-A y cuerpo \{\} en POST /\allowbreak{}api/\allowbreak{}v1/\allowbreak{}routes/\allowbreak{}:\allowbreak{}id/\allowbreak{}complete sobre ruta con paradas resueltas.}{1) Ejecutar cierre de ruta como DRIVER-A. 2) Consultar ruta, parada FAILED y disponibilidad del conductor.}{Ruta pasa a COMPLETED; la parada FAILED sigue fallida y el conductor aún no queda disponible para otra ruta.}{completedAt registrado, pedido FAILED sin marcar DELIVERED y disponibilidad bloqueada hasta retorno.}
\caso{CP-078}{\textbf{Bloqueo de nueva salida antes de confirmar regreso (returnedAt).} NR-11, NR-20, NR-23; NP-13, NP-14 (C28.3, Alta)}{JWT ADMIN, routeId de nueva ruta DRAFT de hoy y driverId de DRIVER-A con returnedAt IS NULL.}{1) Confirmar ruta previa COMPLETED y returnedAt nulo. 2) Intentar asignar nueva ruta a DRIVER-A en POST /\allowbreak{}routes/\allowbreak{}:\allowbreak{}id/\allowbreak{}assign.}{La asignación de nueva salida se rechaza con HTTP 409 mientras no se confirme el regreso físico a tienda.}{La segunda ruta permanece DRAFT sin conductor y no existe doble reserva activa para DRIVER-A.}
\caso{CP-079}{\textbf{Confirmación de regreso por titular libera y actualiza disponibilidad.} NR-11, NR-20, NR-23; NP-13, NP-14 (C28.4, Alta)}{Token DRIVER-A y cuerpo \{\} en POST /\allowbreak{}api/\allowbreak{}v1/\allowbreak{}routes/\allowbreak{}:\allowbreak{}id/\allowbreak{}return (probando también con ausencia activa para hoy).}{1) Confirmar regreso como DRIVER-A en móvil/\allowbreak{}API. 2) Consultar ruta, returnedAt y disponibilidad efectiva (con/\allowbreak{}sin ausencia).}{Registra returnedAt, libera reserva de ruta y actualiza disponibilidad; si existe ausencia hoy, sigue no disponible.}{returnedAt persistido, reserva liberada y disponibilidad efectiva respeta ausencias vigentes sin habilitar indebidamente.}
\caso{CP-080}{\textbf{Repetición de regreso a tienda sin duplicar ni liberar dos veces.} NR-11, NR-20, NR-23; NP-13, NP-14 (C28.5, Alta)}{Token DRIVER-A y segunda invocación con cuerpo \{\} a POST /\allowbreak{}api/\allowbreak{}v1/\allowbreak{}routes/\allowbreak{}:\allowbreak{}id/\allowbreak{}return sobre la misma ruta.}{1) Repetir POST /\allowbreak{}api/\allowbreak{}v1/\allowbreak{}routes/\allowbreak{}:\allowbreak{}id/\allowbreak{}return con DRIVER-A. 2) Comparar returnedAt, auditoría y estado del conductor.}{No duplica registro de regreso ni ejecuta doble liberación; responde idempotente o conflicto controlado según contrato.}{returnedAt original se mantiene sin cambios y la disponibilidad/\allowbreak{}reservas del conductor no sufren alteración.}
\caso{CP-081}{\textbf{Intento de cierre o regreso por DRIVER-B rechazado sobre ruta ajena.} NR-11, NR-20, NR-23; NP-13, NP-14 (C28.6, Alta)}{Token de DRIVER-B hacia POST /\allowbreak{}api/\allowbreak{}v1/\allowbreak{}routes/\allowbreak{}:\allowbreak{}id/\allowbreak{}complete y POST /\allowbreak{}api/\allowbreak{}v1/\allowbreak{}routes/\allowbreak{}:\allowbreak{}id/\allowbreak{}return de DRIVER-A.}{1) Invocar /\allowbreak{}complete y /\allowbreak{}return de ruta de DRIVER-A con JWT de DRIVER-B. 2) Consultar ruta y disponibilidad de DRIVER-A.}{Ambas acciones son rechazadas con HTTP 403 o 404 sin completar ni registrar regreso en la ruta ajena.}{Estado, completedAt, returnedAt y reserva de DRIVER-A permanecen exactamente iguales.}
\end{disenio}

\clearpage
% CALIDAD - diseño de 26 casos (CP-082..CP-107)
% GUÍA EQUIPO: Responsable asignado: José Mateo Auz (Calidad).
\paragraph{Módulo Calidad --- Ruteo, incidencias, concurrencia y rendimiento (26 casos: CP-082 a CP-107).}

\begin{disenio}
\caso{CP-082}{\textbf{Agrupación horaria 10:\allowbreak{}00+10:\allowbreak{}20 y separación 15:\allowbreak{}00.} NR-17, NR-18, NR-19; NP-07, NP-08 (C12.1, Alta)}{Pedidos de hoy: P1 (10:\allowbreak{}00), P2 (10:\allowbreak{}20) distantes y P3 (15:\allowbreak{}00) cercano a P1; margen configurado 30 min.}{1) Enviar GET /\allowbreak{}api/\allowbreak{}v1/\allowbreak{}route-proposals con rol ADMIN. 2) Revisar grupos propuestos y orden de paradas.}{Propuesta agrupa P1+P2 (10:\allowbreak{}00+10:\allowbreak{}20) en una ruta y separa P3 (15:\allowbreak{}00) en otra ruta sin mezclarlos por cercanía.}{P1 y P2 comparten ruta por ventana horaria compatible; P3 queda separado pese a menor distancia a P1.}
\caso{CP-083}{\textbf{Ruta válida con 1 parada sin exigir mínimo de 4.} NR-17, NR-18, NR-19; NP-07, NP-08 (C12.2, Alta)}{1 único pedido elegible de hoy con coordenadas viales válidas dentro de cobertura.}{1) Consultar GET /\allowbreak{}api/\allowbreak{}v1/\allowbreak{}route-proposals como ADMIN. 2) Verificar cantidad de rutas y paradas.}{Se genera 1 propuesta de ruta válida con 1 parada, origen en tienda y retorno, sin exigir mínimo de 4 pedidos.}{HTTP 200 con ruta propuesta de 1 parada utilizable y sin error de capacidad mínima.}
\caso{CP-084}{\textbf{Lote de 4 pedidos compatibles forma ruta de 4 paradas.} NR-17, NR-18, NR-19; NP-07, NP-08 (C12.3, Alta)}{4 pedidos elegibles de hoy en misma ventana horaria (ej. 10:\allowbreak{}00-10:\allowbreak{}25) y margen 30 min.}{1) Solicitar GET /\allowbreak{}api/\allowbreak{}v1/\allowbreak{}route-proposals como ADMIN. 2) Verificar número de rutas y paradas.}{El lote se agrupa en 1 ruta máxima de exactamente 4 paradas ordenadas cronológicamente.}{Total de paradas por ruta \ensuremath{\leq} 4; los 4 pedidos compatibles quedan incluidos en una sola ruta.}
\caso{CP-085}{\textbf{Lote de 5 pedidos dividido en máximo 4 paradas por ruta.} NR-17, NR-18, NR-19; NP-07, NP-08 (C12.4, Alta)}{5 pedidos elegibles de hoy en ventana compatible (margen 30 min) con coordenadas viales válidas.}{1) Invocar GET /\allowbreak{}api/\allowbreak{}v1/\allowbreak{}route-proposals como ADMIN. 2) Inspeccionar partición de rutas y paradas.}{El sistema divide los 5 pedidos en 2 rutas (4+1 o 3+2 según métricas viales), nunca una ruta de 5.}{Todas las rutas tienen entre 1 y 4 paradas; suma total de paradas igual a 5 sin duplicados.}
\caso{CP-086}{\textbf{Declarar emergencia pasa ruta a INTERRUPTED y bloquea.} NR-25; NP-11 (C19.1, Alta)}{Ruta IN\_PROGRESS con 1 entrega DELIVERED y 2 paradas pendientes; EMERGENCY\_DECLARE con routeId y requestId.}{1) Ejecutar POST /\allowbreak{}api/\allowbreak{}v1/\allowbreak{}dispatch/\allowbreak{}interventions con EMERGENCY\_DECLARE. 2) Consultar ruta, RouteUnitRecovery y pedidos.}{Ruta pasa a INTERRUPTED, conserva entrega realizada, crea RouteUnitRecovery por unidad pendiente y las bloquea.}{Estado INTERRUPTED, entrega previa intacta, un registro RouteUnitRecovery por unidad pendiente y bloqueo activo.}
\caso{CP-087}{\textbf{Retorno individual desbloquea solo unidad confirmada.} NR-25; NP-11 (C19.2, Alta)}{Dos unidades pendientes (U1, U2) en RouteUnitRecovery; Product.\allowbreak{}stock inicial S; UNIT\_RECOVERY\_CONFIRM para U1.}{1) Enviar UNIT\_RECOVERY\_CONFIRM con recoveryId de U1. 2) Consultar U1, U2 y Product.\allowbreak{}stock.}{U1 registra returnedAt y se desbloquea; U2 permanece bloqueada y Product.\allowbreak{}stock sigue igual a S.}{Solo U1 queda liberada; U2 conserva returnedAt NULL y no se genera movimiento de suma en Product.\allowbreak{}stock.}
\caso{CP-088}{\textbf{Cierre de emergencia con pendientes mantiene bloqueo 409.} NR-26; NP-11, NP-17 (C20.1, Alta)}{Emergencia con U1 sin retornar; acciones EMERGENCY\_CLOSE y luego intento de gestión o reencolado.}{1) Ejecutar EMERGENCY\_CLOSE. 2) Consultar pendientes visibles. 3) Intentar gestión/\allowbreak{}requeue de pedido bloqueado.}{Incidencia cierra dejando unidades sin retorno visibles y pedidos afectados bloqueados con 409 RECOVERY\_PENDING.}{Pendientes siguen listados tras EMERGENCY\_CLOSE y cualquier intento de requeue/\allowbreak{}gestión retorna 409 RECOVERY\_PENDING.}
\caso{CP-089}{\textbf{Cierre de emergencia no libera repartidor si returnedAt NULL.} NR-26; NP-11, NP-17 (C20.2, Alta)}{Ruta en emergencia cerrada con EMERGENCY\_CLOSE mientras existe RouteUnitRecovery con returnedAt IS NULL.}{1) Ejecutar EMERGENCY\_CLOSE. 2) Consultar disponibilidad de REPARTIDOR-A y reservas activas.}{El cierre EMERGENCY\_CLOSE no restablece por sí solo la disponibilidad ni libera al repartidor mientras returnedAt sea NULL.}{DriverReservation permanece activa tras EMERGENCY\_CLOSE y REPARTIDOR-A no aparece disponible para nueva asignación.}
\caso{CP-090}{\textbf{Liberación de repartidor por ADMIN respeta ausencias y 409.} NR-25, NR-26; NP-11, NP-17 (C20.3, Alta)}{Ruta INTERRUPTED con EMERGENCY\_DRIVER\_RELEASE; variante con ausencia activa; segunda llamada repetida.}{1) Ejecutar EMERGENCY\_DRIVER\_RELEASE como ADMIN. 2) Verificar DriverReservation y registro de ausencia. 3) Repetir acción.}{Libera DriverReservation respetando ausencias vigentes (sigue no disponible si ausente); al repetir rechaza 409 DRIVER\_ALREADY\_RELEASED.}{Reserva liberada una sola vez; la ausencia vigente prevalece en disponibilidad y repetición devuelve 409 DRIVER\_ALREADY\_RELEASED.}
\caso{CP-091}{\textbf{Declarar pérdida con stock \ensuremath{\geq} 1 prepara reemplazo y descuenta.} NR-27; NP-12, NP-18 (C21.1, Alta)}{Original no recuperado de ítem con Product.\allowbreak{}stock = S (\ensuremath{\geq} 1); LOSS\_DECLARE con orderId, orderItemId y requestId.}{1) Ejecutar LOSS\_DECLARE como ADMIN. 2) Consultar OrderItemLoss, Product.\allowbreak{}stock y movimientos de inventario.}{Se crea OrderItemLoss con reemplazo PREPARED, Product.\allowbreak{}stock pasa a S-1 y registra movimiento REPLACEMENT\_OUT.}{Estado PREPARED sin esperar retorno original; exactamente 1 unidad descontada (S-1) con trazabilidad REPLACEMENT\_OUT.}
\caso{CP-092}{\textbf{Reemplazo mantiene mismo pedido y no genera nuevo cobro.} NR-27; NP-12, NP-18 (C21.2, Alta)}{Pedido original con 1 PurchaseReceipt y saldo inicial de CLIENTE-A tras compra original.}{1) Consultar pedido, historial de pérdida/\allowbreak{}reemplazo, PurchaseReceipt y saldo tras LOSS\_DECLARE.}{Conserva el mismo pedido e historial de pérdida/\allowbreak{}reemplazo; no crea nuevo PurchaseReceipt ni nuevo débito.}{Conteo de PurchaseReceipt igual a 1, saldo sin cargos adicionales y mismo ID de pedido.}
\caso{CP-093}{\textbf{Declarar pérdida con stock 0 crea PENDING sin stock negativo.} NR-27; NP-18 (C22.1, Alta)}{Ítem con Product.\allowbreak{}stock === 0; solicitud LOSS\_DECLARE válida con orderId, orderItemId y requestId.}{1) Ejecutar LOSS\_DECLARE como ADMIN. 2) Consultar estado de OrderItemLoss, Product.\allowbreak{}stock y movimientos.}{Crea pérdida en estado PENDING sin descuento de stock (se mantiene en 0) ni reemplazo ficticio.}{Estado PENDING, Product.\allowbreak{}stock === 0 (sin valores negativos) y cero movimientos REPLACEMENT\_OUT.}
\caso{CP-094}{\textbf{Contacto y reprogramación tras reemplazo sin entrega hoy.} NR-27; NP-18 (C22.2, Alta)}{Pedido con reemplazo PREPARED sin entrega posible hoy; acciones CONTACT y RESCHEDULE con fecha acordada.}{1) Registrar CONTACT como ADMIN. 2) Ejecutar RESCHEDULE con fecha futura válida. 3) Consultar auditoría y pedido.}{Registra acciones CONTACT y RESCHEDULE en auditoría y actualiza fecha acordada del pedido sin crear entrega ficticia.}{Fecha acordada persistida, historial auditado de CONTACT/\allowbreak{}RESCHEDULE y sin marcas de entrega adelantada.}
\caso{CP-095}{\textbf{Preparación diferida al reponer stock y 409 NO\_FREE\_STOCK.} NR-27; NP-18 (C22.3, Alta)}{Pérdida PENDING con lossId; primero stock = 0, luego reposición a stock = S (\ensuremath{\geq} 1); acción REPLACEMENT\_PREPARE.}{1) Invocar REPLACEMENT\_PREPARE con stock 0. 2) Reponer stock a S. 3) Repetir REPLACEMENT\_PREPARE y validar.}{Con stock 0 rechaza 409 NO\_FREE\_STOCK; tras reponer a S convierte pérdida a PREPARED y descuenta 1 unidad (S-1).}{409 NO\_FREE\_STOCK sin descuento en paso 1; transición PENDING\ensuremath{\rightarrow}PREPARED y stock S-1 en paso 3.}
\caso{CP-096}{\textbf{Recuperación tardía SELLABLE suma +1 sin doble entrega.} NR-28; NP-12, NP-18 (C23.1, Alta)}{Pérdida con lossId y Product.\allowbreak{}stock = S; acción ORIGINAL\_RECOVER con recoveryCondition = SELLABLE.}{1) Ejecutar ORIGINAL\_RECOVER SELLABLE como ADMIN. 2) Consultar Product.\allowbreak{}stock, movimiento RECOVERY\_IN y pedido.}{Product.\allowbreak{}stock aumenta a S+1 con movimiento RECOVERY\_IN, conserva historial de pérdida/\allowbreak{}reemplazo y no duplica entrega.}{Stock exacto S+1, registro de condición SELLABLE, historial íntegro y pedido sin entrega duplicada.}
\caso{CP-097}{\textbf{Recuperación tardía DAMAGED registra estado sin sumar stock.} NR-28; NP-12, NP-18 (C23.2, Alta)}{Pérdida con lossId y Product.\allowbreak{}stock = S; acción ORIGINAL\_RECOVER con recoveryCondition = DAMAGED.}{1) Ejecutar ORIGINAL\_RECOVER DAMAGED como ADMIN. 2) Consultar condición en OrderItemLoss y Product.\allowbreak{}stock.}{Registra recuperación tardía con condición física DAMAGED sin aumentar Product.\allowbreak{}stock (se mantiene en S).}{Condición DAMAGED persistida en OrderItemLoss, Product.\allowbreak{}stock igual a S y sin movimiento RECOVERY\_IN vendible.}
\caso{CP-098}{\textbf{Repetición de ORIGINAL\_RECOVER rechaza 409 sin doble suma.} NR-28; NP-12, NP-18 (C23.3, Alta)}{Mismo lossId ya recuperado previamente; segundo envío de ORIGINAL\_RECOVER con nuevo requestId.}{1) Enviar ORIGINAL\_RECOVER sobre lossId ya recuperado. 2) Verificar código de respuesta, stock y movimientos.}{Operación rechazada con 409 ORIGINAL\_ALREADY\_RECOVERED sin sumar ni descontar unidades adicionales.}{HTTP 409 ORIGINAL\_ALREADY\_RECOVERED; Product.\allowbreak{}stock e historial de movimientos idénticos a antes del reintento.}
\caso{CP-099}{\textbf{Cancelación LOSS\_CANCEL restituye stock y bloquea si no aplica.} NR-27, NR-28; NP-12, NP-18 (C23.4, Alta)}{Pérdida PREPARED con stock S; pérdida recuperada/\allowbreak{}con reclamo activo o retiro en tienda; LOSS\_CANCEL.}{1) Ejecutar LOSS\_CANCEL sobre pérdida PREPARED. 2) Verificar soft-cancel y stock S+1. 3) Probar casos bloqueados.}{Soft-cancel de OrderItemLoss y restituye +1 (REPLACEMENT\_CANCEL\_IN) si estaba PREPARED; rechaza 409 en bloqueados.}{Stock S+1 solo en PREPARED, soft-cancel auditado y HTTP 409 cuando ya hubo recuperación, reclamo activo o retiro.}
\caso{CP-100}{\textbf{Concurrencia UNIT\_RECOVERY\_CONFIRM vs LOSS\_DECLARE en unidad.} NR-25 a NR-29; Transversal (C24.1, Alta)}{Misma unidad U1 pendiente en ruta INTERRUPTED; 2 hilos simultáneos: UNIT\_RECOVERY\_CONFIRM y LOSS\_DECLARE.}{1) Disparar ambas peticiones concurrentemente con requestId distintos. 2) Inspeccionar respuestas y BD.}{Exactamente una operación gana (200) y la otra rechaza con 409 UNIT\_HAS\_ACTIVE\_LOSS o 409 UNIT\_ALREADY\_RECOVERED.}{Sin estado inconsistente ni doble mutación: 1 éxito y 1 conflicto 409 específico según orden de serialización.}
\caso{CP-101}{\textbf{Idempotencia de intervenciones con requestId UUID y conflicto.} NR-25 a NR-29; Transversal (C24.2, Alta)}{Mismo requestId UUID: 1) envío válido, 2) repetición idéntica, 3) reenvío con payload o acción distinta.}{1) Ejecutar intervención válida. 2) Reenviar payload idéntico con mismo requestId. 3) Reenviar payload distinto.}{Repetición idéntica re-entrega auditoría original sin re-ejecutar efecto; payload distinto rechaza 409 ConflictException.}{1 sola mutación de estado, mismo registro de auditoría en replay y 409 ConflictException ante payload alterado.}
\caso{CP-102}{\textbf{Rendimiento generación de rutas con volumen de 100 pedidos.} NR-17 a NR-19; Plan Sec. 4 (C29.1, Complementaria)}{Dataset sintético aislado de 100 pedidos elegibles del día dentro de cobertura vial.}{1) Registrar versión y hardware del entorno. 2) Medir GET /\allowbreak{}api/\allowbreak{}v1/\allowbreak{}route-proposals con 100 pedidos.}{Se registran latencia p50, máximo, tasa de errores, versión y hardware sin clasificar contra un SLA inventado.}{Reporte reproducible con p50, max, errores y especificaciones de entorno/\allowbreak{}versión para línea base de 100 pedidos.}
\caso{CP-103}{\textbf{Rendimiento generación de rutas con volumen de 500 pedidos.} NR-17 a NR-19; Plan Sec. 4 (C29.2, Complementaria)}{Dataset sintético aislado de 500 pedidos elegibles del día con distribución horaria y vial.}{1) Mantener misma versión y hardware que en 100 pedidos. 2) Medir GET /\allowbreak{}api/\allowbreak{}v1/\allowbreak{}route-proposals con 500 pedidos.}{Se documentan p50, máximo, errores o impedimento real observado en entorno aislado sin inventar SLA fijo.}{Evidencia completa de p50, max, errores, versión y hardware para el volumen de 500 pedidos.}
\caso{CP-104}{\textbf{Rendimiento generación de rutas con volumen de 2000 pedidos.} NR-17 a NR-19; Plan Sec. 4 (C29.3, Complementaria)}{Dataset sintético aislado de 2000 pedidos elegibles del día; referencia comparativa 100/\allowbreak{}500.}{1) Medir GET /\allowbreak{}api/\allowbreak{}v1/\allowbreak{}route-proposals con 2000 pedidos. 2) Comparar tiempos y errores frente a 100 y 500 pedidos.}{Se registran p50, máximo, errores, versión y hardware junto con comparativa 100/\allowbreak{}500/\allowbreak{}2000 sin umbral inventado.}{Tabla comparativa reproducible 100/\allowbreak{}500/\allowbreak{}2000 con métricas y entorno reales, sin calificar aprobación por SLA arbitrario.}
\caso{CP-105}{\textbf{Concurrencia Shop última unidad: [201, 409 STOCK\_UNAVAILABLE].} NR-01, NR-05, NR-29; Transversal (C31.1, Alta)}{Producto con stock = 1; 2 solicitudes simultáneas válidas (CLIENTE-A y CLIENTE-B) con tarjetas demo e Idempotency-Key.}{1) Ejecutar ambos POST /\allowbreak{}api/\allowbreak{}v1/\allowbreak{}shop/\allowbreak{}checkout en paralelo. 2) Consultar respuestas, stock, pedidos y recibos.}{Respuestas [201, 409 STOCK\_UNAVAILABLE]; se crea 1 solo pedido y 1 recibo, stock final 0 sin sobreventa.}{Exactamente un 201 y un 409 STOCK\_UNAVAILABLE; stock \ensuremath{\geq} 0 (0 final) y sin débito ni pedido para la rechazada.}
\caso{CP-106}{\textbf{Concurrencia Shop saldo: [201, 422 INSUFFICIENT\_DEMO\_FUNDS].} NR-01, NR-05, NR-29; Transversal (C31.2, Alta)}{Tarjeta DEMO\_LIMITED con saldo 2000c; producto con stock \ensuremath{\geq} 2; 2 compras concurrentes de 1725c.}{1) Lanzar 2 POST /\allowbreak{}api/\allowbreak{}v1/\allowbreak{}shop/\allowbreak{}checkout simultáneos de 1725c. 2) Verificar respuestas, saldo final y stock.}{Respuestas [201, 422 INSUFFICIENT\_DEMO\_FUNDS]; saldo final 275c y stock descontado solo 1 vez.}{1 compra aprobada (201) y 1 rechazada (422 INSUFFICIENT\_DEMO\_FUNDS); saldo final 275c y stock intacto en el rechazo.}
\caso{CP-107}{\textbf{Fallo pre-commit Shop con trigger en purchaseReceipt: rollback.} NR-01, NR-05, NR-29; Transversal (C31.3, Alta)}{Prueba de integración/\allowbreak{}BD aislada: trigger BEFORE INSERT controlado en purchaseReceipt; checkout válido.}{1) Activar trigger controlado en entorno aislado. 2) Ejecutar POST /\allowbreak{}api/\allowbreak{}v1/\allowbreak{}shop/\allowbreak{}checkout. 3) Consultar BD y retirar trigger.}{Checkout falla con HTTP 500 y realiza rollback total: sin pedido, sin PurchaseReceipt, sin descuento de stock ni saldo.}{HTTP 500; conteos de Order, PurchaseReceipt, Product.\allowbreak{}stock y saldo demo idénticos al estado previo a la petición.}
\end{disenio}


% ===========================================================================
\clearpage
\subsubsection{Preparación del entorno}
% GUÍA PARA EL EQUIPO:
% - Responsable del entorno común: Backend (Saimol Jimenez).
% - Antes de compilar la versión final para entregar, completar cada fila de la tabla siguiente
%   con el commit real (git rev-parse --short HEAD), versión de herramienta y observaciones.
Las pruebas se ejecutaron sobre un entorno de validación aislado con datos sintéticos e integrado
por los siguientes componentes:

\begin{enumerate}
    \item API Backend: PostgreSQL 16 aislado, migraciones Prisma 8 y datos sintéticos. Las pruebas de planificación utilizaron un servicio vial simulado; adicionalmente se verificó una instancia real de OSRM.
    \item Aplicación Web: compilación y ejecución con \texttt{pnpm install} y \texttt{pnpm dev} conectada a la API local aislada.
    \item Aplicación Móvil: cliente \texttt{app-movil-logistica} en Flutter (\texttt{flutter pub get}) ejecutado sobre emulador o dispositivo Android.
    \item Carga y Concurrencia: JMeter con lotes sintéticos de \texttt{100}, \texttt{500} y \texttt{2000} pedidos, además de entorno de base de datos aislado para validación de rollback (\texttt{CP-107}).
\end{enumerate}

\begin{center}
\footnotesize
\begin{tabularx}{\textwidth}{|l|P{3.2cm}|P{2.8cm}|L|}
\hline
\textbf{Componente} & \textbf{Commit (git rev-parse --short HEAD)} & \textbf{Versión de herramienta} & \textbf{Observaciones} \\ \hline
Backend & \texttt{cac8fec} & Node: v26.9.0 & Pruebas locales adicionales no versionadas; PostgreSQL temporal y datos sintéticos. \\ \hline
Frontend web & No registrado & Node: no registrado & Aplicación local conectada a API y base de datos aisladas; pruebas web y de API. \\ \hline
App móvil & \texttt{89f7050} (backend \texttt{9b17cce}) & Flutter: 3.41.3; Node: v26.11.1 & Dispositivo físico Xiaomi M2101K6G, Android 13 (API 33), por USB con \texttt{adb reverse tcp:3003}; base aislada \texttt{logistica\_ape4\_movil}. \\ \hline
Base de datos & No aplica & PostgreSQL: 16 & Bases aisladas y datos sintéticos para las ejecuciones registradas. \\ \hline
Motor vial y rendimiento & Backend \texttt{7eb9a5a}; corrección PR \#76 & OSRM: v5.26.0 (mld, Ecuador); JMeter: 5.6.3 (Java 25) & Intel Core i9-14900HX (32 hilos), 31,7 GB; servicios y generador en la misma laptop. La medición corregida se registró sobre \texttt{7eb9a5a} con el arreglo posteriormente incluido en \texttt{9b8557c}. \\ \hline
\end{tabularx}
\end{center}

% ===========================================================================
\subsubsection{Ejecución de pruebas}
\setcounter{secnumdepth}{4}
\titleformat{\paragraph}[block]{\normalfont\small\bfseries}{\theparagraph}{0.6em}{}
\titlespacing*{\paragraph}{0pt}{8pt}{3pt}
\setcounter{paragraph}{0}
% GUÍA PARA EL EQUIPO (instrucciones internas; no se muestran en el PDF):
% Cada integrante completa únicamente su archivo secciones/<rol>/ejecucion.tex:
% - Completar en cada \fila{ID}{resultado obtenido}{estado}{evidencia}{defecto}:
%   * Estado: PASS (cumple criterio), FAIL (no cumple; registrar defecto), BLOCKED (bloqueado; explicar) o N/A (fuera de alcance; justificar).
%   * Capturas: guardar cada captura representativa en evidencias/CP-XXX.png (el documento la inserta automáticamente).
%   * API/automatizadas: indicar el log o reporte real (Jest, Newman, Playwright o JMeter).
% - Ningún caso debe marcarse PASS sin haber sido ejecutado realmente.
A continuación se presentan las fichas de ejecución y sus evidencias. En cada una se consigna el
resultado registrado; las pruebas con reejecuciones conservan sus antecedentes cuando corresponde.

\clearpage
% BACKEND - fichas individuales de ejecución (27 casos: CP-001..CP-027)
% GUÍA EQUIPO (Saimol Jimenez): completar los parámetros de \casoejec{ID}{título}{resultado obtenido}{PASS|FAIL|BLOCKED|N/A}{evidencia o log}{DEF-XX}.
% - Estado: solo PASS, FAIL, BLOCKED o N/A.
% - Defecto: ID del defecto (ej. DEF-01) si falló, o Ninguno si pasó.
% - Resultado obtenido: describir lo observado al ejecutar la prueba.
% - Si existe evidencias/<ID>.png, la imagen se inserta automáticamente dentro de su ficha.
\par\medskip\noindent\textbf{Módulo Backend y reglas de negocio (27 casos: CP-001 a CP-027).}\par\smallskip
Se ejecutaron pruebas automatizadas de la API REST en un entorno aislado. Las imágenes son representaciones de los resúmenes conjuntos reales de Jest, no capturas individuales de las respuestas HTTP; el código de prueba identifica las aserciones de cada caso. Las pruebas de ruteo emplearon un servicio vial simulado, no OSRM real.

\casoejec{CP-001}{Compra ADMIN para cliente con destinatario de regalo}{Compra registrada con destinatario independiente; lecturas ADMIN y cliente correctas y stock descontado.}{PASS}{Suite automatizada CP-001 a CP-012; ver log backend comercial.}{Ninguno}
\casoejec{CP-002}{Rechazo compra ADMIN cliente sin teléfono de comprador}{HTTP 400; pedidos, comprobantes y stock permanecieron sin cambios.}{PASS}{Suite automatizada CP-001 a CP-012; ver log backend comercial.}{Ninguno}
\casoejec{CP-003}{Compra ADMIN invitado y lectura de comprobante}{Compra invitado exitosa, comprobante comercial consultado y ningún usuario nuevo creado.}{PASS}{Suite automatizada CP-001 a CP-012; ver log backend comercial.}{Ninguno}
\casoejec{CP-004}{Rechazo compra ADMIN invitado sin teléfono obligatorio}{Teléfono ausente y vacío rechazados con HTTP 400, sin cambios en pedidos, comprobantes, usuarios ni stock.}{PASS}{Suite automatizada CP-001 a CP-012; ver log backend comercial.}{Ninguno}
\casoejec{CP-005}{Compra invitado con teléfono de cliente sin vínculo}{La compra invitada no se vinculó al cliente con el mismo teléfono; su historial no la mostró y el detalle devolvió 404.}{PASS}{Suite automatizada CP-001 a CP-012; ver log backend comercial.}{Ninguno}
\casoejec{CP-006}{Pago externo declarado ADMIN en efectivo}{Pago externo en efectivo declarado y comprobante verificados; sin débito de tarjeta simulada.}{PASS}{Suite automatizada CP-001 a CP-012; ver log backend comercial.}{Ninguno}
\casoejec{CP-007}{Pago externo declarado ADMIN por transferencia}{Transferencia declarada y comprobante registrados sin referencia bancaria ni intento de pago simulado.}{PASS}{Suite automatizada CP-001 a CP-012; ver log backend comercial.}{Ninguno}
\casoejec{CP-008}{Cambio de tasa IVA tras cotizar (409 SHOP\_SETTINGS\_CHANGED)}{Tras cambiar el IVA se obtuvo HTTP 409 SHOP\_SETTINGS\_CHANGED sin cambios en pedido, comprobante, stock ni saldo.}{PASS}{Suite automatizada CP-001 a CP-012; ver log backend comercial.}{Ninguno}
\casoejec{CP-009}{Cambio de precio tras cotizar e inmutabilidad 15\% IVA}{Precio modificado rechazado con HTTP 409; tras recotizar, el comprobante conservó 7300/1095/8395 centavos ante cambios posteriores.}{PASS}{Suite automatizada CP-001 a CP-012; ver log backend comercial.}{Ninguno}
\casoejec{CP-010}{Corrección autorizada de contacto ADMIN y auditoría}{ADMIN actualizó ambos teléfonos; la revisión pasó de 0 a 1 y se registró el historial auditado.}{PASS}{Suite automatizada CP-001 a CP-012; ver log backend comercial.}{Ninguno}
\casoejec{CP-011}{Aislamiento de corrección de contacto en recibo y pedido}{La corrección de P1 conservó su comprobante original y no alteró el pedido de control P2.}{PASS}{Suite automatizada CP-001 a CP-012; ver log backend comercial.}{Ninguno}
\casoejec{CP-012}{Rechazo de corrección de contacto por rol no autorizado}{CLIENTE y REPARTIDOR recibieron HTTP 403; pedido e historial permanecieron iguales.}{PASS}{Suite automatizada CP-001 a CP-012; ver log backend comercial.}{Ninguno}
\casoejec{CP-013}{Horario ADMIN límite inicial 07:\allowbreak{}00 America/\allowbreak{}Guayaquil}{HTTP 201 a las 07:00; pedido PENDING y comprobante consultados, con descuento de una unidad de stock.}{PASS}{Suite CP-013 a CP-027; resumen conjunto de Jest.}{Ninguno}
\casoejec{CP-014}{Horario ADMIN límite final 17:\allowbreak{}00 America/\allowbreak{}Guayaquil}{HTTP 201 a las 17:00; pedido PENDING y comprobante consultados, con descuento de una unidad de stock.}{PASS}{Suite CP-013 a CP-027; resumen conjunto de Jest.}{Ninguno}
\casoejec{CP-015}{Rechazo horario ADMIN inválido 06:\allowbreak{}59 (422)}{HTTP 422 DELIVERY\_OUTSIDE\_HOURS; pedidos, comprobantes y stock sin cambios.}{PASS}{Suite CP-013 a CP-027; resumen conjunto de Jest.}{Ninguno}
\casoejec{CP-016}{Rechazo horario ADMIN inválido 17:\allowbreak{}01 (422)}{HTTP 422 DELIVERY\_OUTSIDE\_HOURS; pedidos, comprobantes y stock sin cambios.}{PASS}{Suite CP-013 a CP-027; resumen conjunto de Jest.}{Ninguno}
\casoejec{CP-017}{Agrupación temporal por extremos (10:\allowbreak{}00, 10:\allowbreak{}25, 10:\allowbreak{}50)}{HTTP 200: dos propuestas, una con 10:00 y 10:25 y otra con 10:50; sin persistir ruta. Motor vial simulado.}{PASS}{Suite CP-013 a CP-027; resumen conjunto de Jest.}{Ninguno}
\casoejec{CP-018}{Agrupación por extremos con separación de 29 minutos}{HTTP 200: los dos pedidos se agruparon en una propuesta cronológica con margen de 30 minutos. Motor vial simulado.}{PASS}{Suite CP-013 a CP-027; resumen conjunto de Jest.}{Ninguno}
\casoejec{CP-019}{Agrupación por extremos con separación exacta de 30 min}{HTTP 200: los pedidos separados exactamente 30 minutos se agruparon en una propuesta. Motor vial simulado.}{PASS}{Suite CP-013 a CP-027; resumen conjunto de Jest.}{Ninguno}
\casoejec{CP-020}{Separación de rutas por extremos con 31 minutos}{HTTP 200: los pedidos separados 31 minutos produjeron dos propuestas individuales. Motor vial simulado.}{PASS}{Suite CP-013 a CP-027; resumen conjunto de Jest.}{Ninguno}
\casoejec{CP-021}{Cambio de margen horario (5..180) y desempate OSRM}{Se aceptaron los límites 5 y 180, se rechazaron 4 y 181 y se aplicó margen 60. El agrupamiento y desempate se verificaron con matrices simuladas; adicionalmente, una prueba aislada del motor OSRM real respondió correctamente.}{PASS}{Suites de configuración, ruteo simulado y motor OSRM real; resumen en imagen.}{Ninguno}
\casoejec{CP-022}{Creación de pedido sin fecha programada (scheduledDeliveryDate nulo)}{HTTP 201; pedido PENDING con fecha nula confirmado por consulta y dos unidades descontadas del stock.}{PASS}{Suite CP-013 a CP-027; resumen conjunto de Jest.}{Ninguno}
\casoejec{CP-023}{Exclusión en propuestas y 409 en asignación sin fecha}{La ejecución anterior falló al exigir el pedido en \texttt{excluded}. Tras aclarar la regla, el pedido con fecha nula se identificó en el listado ADMIN, no apareció en propuestas y ambos intentos de asignación respondieron 409 sin crear ruta ni reclamo.}{PASS}{Nueva ejecución aislada: 15/15 pruebas de la suite aprobadas; antecedente DEF-08 en la imagen.}{DEF-08 cerrado}
\casoejec{CP-024}{Idempotencia Shop checkout con replay 201}{Ambas respuestas fueron HTTP 201 con los mismos identificadores; un pedido, un comprobante y un solo descuento y débito.}{PASS}{Suite CP-013 a CP-027; resumen conjunto de Jest.}{Ninguno}
\casoejec{CP-025}{Recuperación Shop ante timeout con mismo Idempotency-Key}{Tras corte de conexión, el reintento respondió 201 y quedó un solo pedido, comprobante, descuento y débito. La suite previa de recuperación también verificó los identificadores recuperados.}{PASS}{Suite CP-013 a CP-027 y suite aislada de recuperación Shop.}{Ninguno}
\casoejec{CP-026}{Conflicto Idempotency-Key con payload distinto (409)}{HTTP 409 IDEMPOTENCY\_KEY\_REUSED; pedidos, comprobantes, stock y saldo sin cambios adicionales.}{PASS}{Suite CP-013 a CP-027; resumen conjunto de Jest.}{Ninguno}
\casoejec{CP-027}{Liberación auditada de DRAFT caducado (RELEASE\_DRAFT)}{La ejecución anterior falló al esperar elegibilidad hoy sin reprogramar. Con la regla aclarada, se verificaron el 409, la liberación y replay auditado, la fecha original conservada y la reprogramación mediante CONTACT y RESCHEDULE. El pedido apareció en propuestas para el día futuro acordado. Motor vial simulado.}{PASS}{Nueva ejecución aislada: 15/15 pruebas de la suite aprobadas; antecedente DEF-09 en la imagen.}{DEF-09 cerrado}

\clearpage
% FRONTEND WEB - fichas individuales de ejecución (27 casos: CP-028..CP-054)
% GUÍA EQUIPO (Kerly Chicaiza): completar los parámetros de \casoejec{ID}{título}{resultado obtenido}{PASS|FAIL|BLOCKED|N/A}{evidencia o log}{DEF-XX}.
% - Estado: solo PASS, FAIL, BLOCKED o N/A.
% - Defecto: ID del defecto (ej. DEF-01) si falló, o Ninguno si pasó.
% - Resultado obtenido: describir lo observado al ejecutar la prueba.
% - Si existe evidencias/<ID>.png, la imagen se inserta automáticamente dentro de su ficha.
\par\medskip\noindent\textbf{Módulo Web --- Tienda en línea y panel de rutas (27 casos: CP-028 a CP-054).}\par\smallskip
Pruebas ejecutadas sobre la aplicación web React conectada al backend aislado y verificadas mediante Playwright y Vitest.

\casoejec{CP-028}{Historial unificado Mis pedidos con compras SHOP y ADMIN propias}{}{}{Captura UI / reporte}{}
\casoejec{CP-029}{Bloqueo en lista y detalle ante pedido ajeno o invitado coincidente}{}{}{Captura UI / reporte}{}
\casoejec{CP-030}{Pedido sin PurchaseReceipt y sin importes inventados}{}{}{Captura UI / reporte}{}
\casoejec{CP-031}{Cobertura tienda: punto interior al polígono aceptado en web}{}{}{Captura UI / reporte}{}
\casoejec{CP-032}{Cobertura tienda: punto exactamente en el borde del polígono}{}{}{Captura UI / reporte}{}
\casoejec{CP-033}{Cobertura tienda: punto exterior rechazado con 422 en web}{}{}{Captura UI / reporte}{}
\casoejec{CP-034}{Cambio de cobertura A a B conserva pedidos ya aceptados}{}{}{Captura UI / reporte}{}
\casoejec{CP-035}{Compra nueva tras cambio de cobertura aplica área B vigente}{}{}{Captura UI / reporte}{}
\casoejec{CP-036}{Horario Shop límite inicial 07:\allowbreak{}00 America/\allowbreak{}Guayaquil aceptado}{}{}{Captura UI / reporte}{}
\casoejec{CP-037}{Horario Shop límite final 17:\allowbreak{}00 America/\allowbreak{}Guayaquil aceptado}{}{}{Captura UI / reporte}{}
\casoejec{CP-038}{Horario Shop inválido 06:\allowbreak{}59 rechazado sin orden ni cobro}{}{}{Captura UI / reporte}{}
\casoejec{CP-039}{Horario Shop inválido 17:\allowbreak{}01 rechazado sin orden ni cobro}{}{}{Captura UI / reporte}{}
\casoejec{CP-040}{Cambio de horario en shop-settings genera 409 al confirmar}{}{}{Captura UI / reporte}{}
\casoejec{CP-041}{Verificación UI web de aviso de hora de entrega aproximada}{}{}{Captura UI / reporte}{}
\casoejec{CP-042}{Propuesta de rutas en panel ADMIN es solo lectura sin persistir}{}{}{Captura UI / reporte}{}
\casoejec{CP-043}{Confirmación de borrador DRAFT en hoy persiste ruta sin conductor}{}{}{Captura UI / reporte}{}
\casoejec{CP-044}{Asignación separada de ruta DRAFT a conductor disponible}{}{}{Captura UI / reporte}{}
\casoejec{CP-045}{Rechazo de asignación a conductor no disponible o ausente (409)}{}{}{Captura UI / reporte}{}
\casoejec{CP-046}{Motor vial OSRM indisponible muestra aviso sin rutas en línea recta}{}{}{Captura UI / reporte}{}
\casoejec{CP-047}{Preservación de pedidos pendientes y rutas tras falla OSRM}{}{}{Captura UI / reporte}{}
\casoejec{CP-048}{Recuperación tras indisponibilidad OSRM genera propuestas válidas}{}{}{Captura UI / reporte}{}
\casoejec{CP-049}{Origen y retorno común en tienda única reflejado en propuesta y mapa}{}{}{Captura UI / reporte}{}
\casoejec{CP-050}{Paradas y trazado vial en panel ADMIN coinciden con motor vial}{}{}{Captura UI / reporte}{}
\casoejec{CP-051}{Distancias y duraciones mostradas como estimadas sin tráfico en vivo}{}{}{Captura UI / reporte}{}
\casoejec{CP-052}{Regresión Shop: compra válida con DEMO\_APPROVED crea pedido y recibo}{}{}{Captura UI / reporte}{}
\casoejec{CP-053}{Regresión Shop: carrito mixto con agotado rechazado sin efectos}{}{}{Captura UI / reporte}{}
\casoejec{CP-054}{Regresión Shop: precio o total manipulado en cliente no es aceptado}{}{}{Captura UI / reporte}{}

\clearpage
% MÓVIL - fichas individuales de ejecución (27 casos: CP-055..CP-081)
% GUÍA EQUIPO (Carlos Ramos): completar los parámetros de \casoejec{ID}{título}{resultado obtenido}{PASS|FAIL|BLOCKED|N/A}{evidencia o log}{DEF-XX}.
% - Estado: solo PASS, FAIL, BLOCKED o N/A.
% - Defecto: ID del defecto (ej. DEF-01) si falló, o Ninguno si pasó.
% - Resultado obtenido: describir lo observado al ejecutar la prueba.
% - Si existe evidencias/<ID>.png, la imagen se inserta automáticamente dentro de su ficha.
\par\medskip\noindent\textbf{Módulo Móvil --- Seguridad y operación de entregas (27 casos: CP-055 a CP-081).}\par\smallskip
En esta sección se presentan los resultados de la ejecución de los 27 casos de prueba del módulo
móvil (\texttt{CP-055} a \texttt{CP-081}): control de acceso por rol y por recurso, planificación del
día operativo, urgencias, entregas fallidas y su gestión posterior, recorrido de extremo a extremo y
cierre de ruta con regreso a tienda.

Todas las pruebas se ejecutaron en local. El backend NestJS (commit \texttt{9b17cce}) se levantó en
el puerto \texttt{3003} sobre una base PostgreSQL~16 aislada creada solo para esta ejecución
(\texttt{logistica\_ape4\_movil}, 22 migraciones aplicadas y datos demo sintéticos), con el simulador
de pagos habilitado. La aplicación Flutter (\texttt{app-movil-logistica}, commit \texttt{89f7050}) se
compiló en modo \emph{debug} y se instaló en un teléfono físico Xiaomi M2101K6G con Android~13
conectado por cable USB; el teléfono alcanzó la API local mediante \texttt{adb reverse tcp:3003}.

Los casos de API se automatizaron con el script \texttt{evidencias/scripts/ape4-movil.mjs}, que inicia
sesión con cuentas reales (\texttt{ADMIN}, \texttt{CLIENT-A}, \texttt{CLIENT-B}, \texttt{DRIVER-A},
\texttt{DRIVER-B}) y llama a la API por HTTP con sus tokens JWT. Cada caso guarda en
\texttt{evidencias/json/CP-XXX.json} las peticiones, las respuestas y cada verificación del criterio de
aceptación; cuando el criterio exige comprobar persistencia (reservas, cuentas, retiros) se consultó
además la base de datos. En los casos de rechazo se comparó el estado antes y después de cada intento:
un caso solo se consideró exitoso si, además del código HTTP esperado, no dejó efectos parciales. Las
ausencias registradas durante las pruebas se anularon al terminar cada caso.

Los casos que involucran la pantalla del repartidor (\texttt{CP-057}, \texttt{CP-060},
\texttt{CP-067}, \texttt{CP-073} a \texttt{CP-077} y \texttt{CP-079}) se repitieron en el teléfono con
el script de apoyo \texttt{evidencias/scripts/ape4-movil-telefono.mjs}, que registra el estado real
del backend después de cada acción en la app (\texttt{evidencias/json/movil-telefono.json}). La
ejecución en el teléfono reveló dos defectos de la app: no carga rutas creadas sin zona
(\texttt{DEF-06}) y no ofrece las acciones de entrega fallida, cierre de ruta ni confirmación de
regreso (\texttt{DEF-07}). Ambos se detallan en las fichas y en el registro de incidentes.

% ---------------------------------------------------------------------------
\par\medskip\noindent\textbf{Seguridad y control de acceso (CP-055 a CP-060).}\par\smallskip
Se comprobó que cada actor solo accede a sus propios recursos y que ningún rechazo filtra datos ni
modifica el estado.

\casoejec{CP-055}{Acceso sin sesión JWT a rutas, pedidos y configuración}{Las tres peticiones sin
cabecera \texttt{Authorization} (\texttt{GET /drivers/me/routes}, \texttt{GET /shop/orders} y
\texttt{GET /admin/shop-settings}) respondieron HTTP 401. Ninguna respuesta incluyó listados, rutas,
paradas ni parámetros de configuración.}{PASS}{Script \texttt{ape4-movil.mjs};
JSON: \texttt{evidencias/json/CP-055.json}}{Ninguno}

\casoejec{CP-056}{Aislamiento entre CLIENT-B y pedido/\allowbreak{}comprobante de CLIENT-A}{Con el JWT de
CLIENT-B sobre una compra Shop de CLIENT-A: \texttt{GET /shop/orders/:id} y \texttt{/receipt}
respondieron 404, y \texttt{GET /orders/:id} y \texttt{PATCH /orders/:id/status} respondieron 403.
Ninguna respuesta contenía el destinatario, la dirección, el número comercial ni el total de CLIENT-A,
y el pedido consultado como ADMIN quedó idéntico antes y después.}{PASS}{JSON:
\texttt{evidencias/json/CP-056.json}}{Ninguno}

\casoejec{CP-057}{Aislamiento de DRIVER-B ante ruta y parada de DRIVER-A}{\emph{API:} Mi Ruta de
DRIVER-B no incluyó la ruta de DRIVER-A; consultarla (\texttt{GET /routes/:id}), iniciarla, registrar
un fallo en su parada y cerrarla como DRIVER-B respondieron 403, y la ruta y sus paradas conservaron el
mismo estado. \emph{Teléfono:} con la sesión del segundo repartidor la app mostró ``No tienes rutas
para hoy'', sin listar la ruta asignada al primer repartidor.}{PASS}{JSON:
\texttt{evidencias/json/CP-057.json}; captura del teléfono}{Ninguno}

\casoejec{CP-058}{Registro público con rol ADMIN rechazado sin escalación}{\texttt{POST
/auth/register} con un correo nuevo y el campo adicional \texttt{role: ADMIN} respondió 400. El inicio
de sesión posterior con esas credenciales respondió 401. En la base de datos no existe ninguna cuenta
con ese correo y el número de cuentas ADMIN siguió siendo 1 (la del seed).}{PASS}{JSON y consulta a
la base: \texttt{evidencias/json/CP-058.json}}{Ninguno}

\casoejec{CP-059}{Intento de CLIENT de modificar configuración o cobertura}{Como CLIENT-A,
\texttt{PATCH /admin/shop-settings} y \texttt{PUT /admin/store-coverage} (nuevo polígono) respondieron
403. La configuración de la tienda y la cobertura (ubicación y polígono) leídas como ADMIN fueron
idénticas antes y después.}{PASS}{JSON comparativo: \texttt{evidencias/json/CP-059.json}}{Ninguno}

\casoejec{CP-060}{Acceso autorizado por titular, ADMIN y DRIVER sobre recursos propios}{CLIENT-A
leyó su pedido y su comprobante (200); su historial devolvió 3 pedidos, todos propios. ADMIN consultó
la configuración y la bandeja de despacho (200). DRIVER-A obtuvo 200 en Mi Ruta solo con su ruta
asignada (0 rutas de otros repartidores). En el teléfono, el repartidor vio su ruta del día (zona
Centro, Asignada, 0 de 2 entregados); esta prueba se hizo con una ruta creada con zona, por el defecto
\texttt{DEF-06} descrito en \texttt{CP-073}.}{PASS}{JSON: \texttt{evidencias/json/CP-060.json}; captura
del teléfono}{Ninguno}

% ---------------------------------------------------------------------------
\par\medskip\noindent\textbf{Planificación del día operativo y urgencias (CP-061 a CP-066).}\par\smallskip
Se verificó que solo los pedidos del día operativo (America/Guayaquil) se conviertan en rutas y que
las urgencias nunca alteren una ruta iniciada sin decisión del ADMIN.

\casoejec{CP-061}{Pedido de hoy elegible en propuesta, ruta y asignación}{Con dos pedidos para hoy
(16:00 y 16:10), \texttt{GET /route-proposals} respondió 200 con \texttt{isToday=true} e incluyó
ambos. \texttt{POST /routes} creó la ruta en \texttt{DRAFT} (201) y \texttt{POST
/routes/:id/assign} la dejó \texttt{ASSIGNED} a DRIVER-A (200), con 2 paradas y la fecha de
hoy.}{PASS}{JSON: \texttt{evidencias/json/CP-061.json}}{Ninguno}

\casoejec{CP-062}{Pedido futuro admitido en compra y bloqueado en ruta/\allowbreak{}assign}{La compra
Shop y el pedido ADMIN para mañana se crearon (201). La vista previa de mañana respondió 200 con
\texttt{isToday=false} y \texttt{canConfirm=false} e incluyó ambos pedidos. \texttt{POST /routes} para
mañana respondió 409 \texttt{ROUTE\_DATE\_NOT\_TODAY} y la asignación directa \texttt{POST
/routes/assign} también 409. No se creó ninguna ruta y los dos pedidos siguieron
\texttt{PENDING}.}{PASS}{JSON: \texttt{evidencias/json/CP-062.json}}{Ninguno}

\casoejec{CP-063}{Ausencia registrada hoy impide asignar repartidor}{Se registró una ausencia de hoy
(201) a un repartidor sin rutas, para que la ausencia fuera la única causa posible. Su disponibilidad
devolvió \texttt{isAvailable=false} con la ausencia indicada y el intento de asignarle una ruta
respondió 409. La ruta quedó en \texttt{DRAFT} sin conductor y no se creó ninguna reserva (0
reservas).}{PASS}{JSON: \texttt{evidencias/json/CP-063.json}}{Ninguno}

\casoejec{CP-064}{Ausencia futura no altera ni reasigna rutas activas de hoy}{Con la ruta de hoy de
DRIVER-A asignada, se registró una ausencia de mañana a pasado mañana (201, \texttt{ACTIVE}). La ruta
mantuvo el mismo conductor, estado, paradas e intentos y su reserva siguió activa (1 reserva): no hubo
reasignación, cancelación ni liberación automática. La ausencia se anuló al terminar.}{PASS}{JSON
comparativo: \texttt{evidencias/json/CP-064.json}}{Ninguno}

\casoejec{CP-065}{Pedido urgente con ruta IN\_PROGRESS no altera ruta iniciada}{Con la ruta de
DRIVER-A ya iniciada (\texttt{IN\_PROGRESS}), un pedido urgente creado para hoy obtuvo la propuesta
\texttt{NEW\_ROUTE} en \texttt{GET /dispatch/orders/:id/proposal}, sin proponer la ruta iniciada. Las
paradas, el orden y el estado de la ruta iniciada fueron idénticos antes y después.}{PASS}{JSON
comparativo: \texttt{evidencias/json/CP-065.json}}{Ninguno}

\casoejec{CP-066}{Decisión manual ADMIN en urgencias: INSERT con cupo o NEW\_ROUTE}{Mientras había
rutas no iniciadas con cupo, cada urgente obtuvo la propuesta \texttt{INSERT}. La ruta destino no
cambió hasta que el ADMIN confirmó explícitamente (\texttt{CONFIRM}, 200 \texttt{APPLIED}); cada
confirmación agregó exactamente una parada con secuencia contigua. Se hicieron 5 inserciones en dos
rutas hasta llenarlas (4 paradas cada una). El siguiente urgente obtuvo \texttt{NEW\_ROUTE} y, al
confirmarlo con un repartidor libre, se creó una ruta nueva. Ninguna ruta superó 4 paradas, las
decisiones quedaron auditadas con el actor ADMIN y la ruta iniciada no recibió
inserciones.}{PASS}{JSON y auditoría de despacho: \texttt{evidencias/json/CP-066.json}}{Ninguno}

% ---------------------------------------------------------------------------
\par\medskip\noindent\textbf{Entrega fallida y gestión posterior (CP-067 a CP-071).}\par\smallskip
Se verificó que una visita fallida se conserve en el historial, que solo el repartidor titular pueda
registrarla y que las gestiones del ADMIN (reprogramar o retiro en tienda) no dupliquen entregas.

\casoejec{CP-067}{Entrega fallida en móvil conserva parada e historial pendiente}{\emph{API:} DRIVER-A
registró la entrega fallida con motivo \texttt{RECIPIENT\_ABSENT} (200). La parada pasó a
\texttt{FAILED} conservando el intento con su motivo, el pedido no quedó \texttt{DELIVERED} ni
\texttt{CANCELLED} y el seguimiento del cliente lo mostró pendiente de gestión.
\emph{Teléfono:} la app no permite registrar el fallo; la parada en curso solo ofrece ``En camino'' y
``Entregado''. El fallo tuvo que registrarse por API y la app solo lo mostró después (``Intento
fallido''). El paso ``registrar entrega fallida desde móvil'' no se puede ejecutar.}{FAIL}{JSON:
\texttt{evidencias/json/CP-067.json} y \texttt{movil-telefono.json}; captura del teléfono}{DEF-07}

\casoejec{CP-068}{Registro de fallo por actor o conductor no autorizado rechazado}{El intento de
registrar un fallo sobre la parada de DRIVER-A respondió 403 tanto con el JWT de DRIVER-B como con el
de CLIENT-A. La parada siguió \texttt{PENDING} sin intentos y el pedido no cambió de
estado.}{PASS}{JSON: \texttt{evidencias/json/CP-068.json}}{Ninguno}

\casoejec{CP-069}{Gestión post-fallo ADMIN: CONTACT y RESCHEDULE con nueva fecha}{Tras una visita
fallida y la ruta cerrada, \texttt{CONTACT} con acuerdo explícito (mañana 12:00, ventana 11:30--12:30)
y \texttt{RESCHEDULE} con ese acuerdo respondieron 200 \texttt{APPLIED}. El pedido quedó
\texttt{PENDING}, sin ruta y con la nueva fecha (\texttt{2026-11-15T17:00:00Z}). La ruta original
conservó la parada \texttt{FAILED} con su intento, no quedó ninguna asignación activa duplicada y el
pedido apareció en la propuesta de mañana.}{PASS}{JSON: \texttt{evidencias/json/CP-069.json}}{Ninguno}

\casoejec{CP-070}{Gestión post-fallo ADMIN: acuerdo de retiro en tienda (PICKUP\_AGREE)}{\texttt{PICKUP\_AGREE}
respondió 200 \texttt{APPLIED}. El pedido quedó \texttt{PENDING} con el retiro en
\texttt{PENDING\_PICKUP} (sin \texttt{pickedUpAt}), sin marcarse entregado. El seguimiento del cliente
mostró el hito ``Pendiente de retiro en tienda'', la visita fallida se conservó y el stock no cambió
(104 $\rightarrow$ 104).}{PASS}{JSON: \texttt{evidencias/json/CP-070.json}}{Ninguno}

\casoejec{CP-071}{Confirmación de retiro (PICKUP\_CONFIRM) y rechazo de repetición}{La primera
\texttt{PICKUP\_CONFIRM} respondió 200: el pedido quedó \texttt{DELIVERED} con el retiro en
\texttt{PICKED\_UP}. La repetición se rechazó con 409, pero con el mensaje genérico ``Commercially
resolved order cannot be managed for pickup'' y sin el código \texttt{PICKUP\_ALREADY\_CONFIRMED} que
exige el caso. La integridad sí se cumplió: un solo registro de retiro y stock sin cambio (104
$\rightarrow$ 104).}{FAIL}{JSON: \texttt{evidencias/json/CP-071.json}}{DEF-05}

% ---------------------------------------------------------------------------
\par\medskip\noindent\textbf{Recorrido de extremo a extremo (CP-072 a CP-075).}\par\smallskip
Se recorrió el flujo completo, desde la compra del cliente hasta la entrega en el teléfono del
repartidor, contrastando lo que muestra la app con lo que registra el servidor.

\casoejec{CP-072}{Recorrido E2E: creación de compra y verificación de estado inicial}{CLIENT-A compró 2
unidades con la tarjeta demo aprobada (201). El pedido se creó en \texttt{PENDING} con importes
congelados: subtotal 5180, IVA 15\,\% 777 y total 5957 centavos, igual a lo cotizado. El seguimiento y
el comprobante de CLIENT-A mostraron el mismo pedido y total, y el stock bajó de 109 a 107.}{PASS}{JSON:
\texttt{evidencias/json/CP-072.json}}{Ninguno}

\casoejec{CP-073}{Recorrido E2E: propuesta, confirmación y asignación en día operativo}{\emph{API:} la
propuesta de hoy incluyó el pedido; el borrador se persistió (201, \texttt{DRAFT}) y quedó
\texttt{ASSIGNED} a DRIVER-A, con una sola reserva y una sola aparición en \texttt{GET
/drivers/me/routes}. \emph{Teléfono:} la ruta confirmada desde la propuesta no tiene zona
(\texttt{"zone": null}) y la app no la muestra. Mi Ruta respondió ``No se pudieron cargar tus rutas. El
servidor devolvió rutas inválidas''. La app exige que \texttt{zone} sea un objeto, y con una sola ruta
sin zona el repartidor no ve ninguna de sus rutas. Con una ruta creada con zona, la app la listó
correctamente.}{FAIL}{JSON: \texttt{evidencias/json/CP-073.json}; captura del
teléfono}{DEF-06}

\casoejec{CP-074}{Recorrido E2E: progreso móvil reflejado en seguimiento del cliente}{\emph{API:} con
la compra de CP-072, cada transición (inicio, \texttt{IN\_TRANSIT}, \texttt{DELIVERED}) se reflejó
igual en el backend y en el seguimiento de CLIENT-A, sin coordenadas GPS en la línea de tiempo.
\emph{Teléfono:} el repartidor inició la ruta y marcó ``En camino'' y ``Entregado'' en la parada de una
compra Shop. Tras cada acción, el seguimiento del cliente coincidió con el backend: \texttt{IN\_TRANSIT}
y luego \texttt{DELIVERED}, con la línea de tiempo PENDING $\rightarrow$ ASSIGNED $\rightarrow$
IN\_TRANSIT $\rightarrow$ DELIVERED.}{PASS}{JSON: \texttt{evidencias/json/CP-074.json} y
\texttt{movil-telefono.json}; captura del teléfono}{Ninguno}

\casoejec{CP-075}{Recorrido E2E: desconexión móvil y recuperación al reconectar}{En el teléfono se
cortó la conexión con el servidor (se quitó \texttt{adb reverse}) y se pulsó ``Iniciar ruta''. La app
no fingió éxito: mostró ``No se pudo conectar con el servidor. Actualiza la lista antes de volver a
intentar iniciar una ruta'', la ruta siguió ``Asignada'' y en el backend continuó \texttt{ASSIGNED} con
\texttt{startedAt} nulo. Al reconectar y actualizar, la app mostró el estado real del servidor, sin
cola offline. El inicio posterior se registró una sola vez.}{PASS}{JSON:
\texttt{evidencias/json/CP-075.json}; captura del teléfono}{Ninguno}

% ---------------------------------------------------------------------------
\par\medskip\noindent\textbf{Cierre de ruta y regreso a tienda (CP-076 a CP-081).}\par\smallskip
Se verificó que la ruta solo se cierre con todas las visitas resueltas y que el repartidor quede
ocupado hasta confirmar su regreso físico a la tienda.

\casoejec{CP-076}{Cierre de ruta móvil rechazado con visita aún PENDING}{Con una parada todavía
\texttt{PENDING}, \texttt{POST /routes/:id/complete} de DRIVER-A respondió 409. La ruta siguió
\texttt{IN\_PROGRESS} sin \texttt{completedAt} y el repartidor siguió ocupado. En el teléfono se
repitió el escenario (parada 1 con intento fallido y parada 2 pendiente): el backend rechazó el cierre
con 409 (``One to four resolved verified visits required'') y la app mostró la ruta ``En curso'' con la
parada pendiente. La app no tiene botón de cierre (\texttt{DEF-07}).}{PASS}{JSON:
\texttt{evidencias/json/CP-076.json}; captura del teléfono}{Ninguno}

\casoejec{CP-077}{Cierre de ruta COMPLETED con visitas resueltas y una FAILED}{DRIVER-A entregó la
segunda parada y cerró la ruta: 200 \texttt{COMPLETED} con \texttt{completedAt}. La parada fallida
siguió \texttt{FAILED}, su pedido no se marcó \texttt{DELIVERED} y el repartidor siguió no disponible
con la ruta como bloqueante hasta su regreso. En el teléfono, después del cierre (hecho por API con el
token del repartidor, ver \texttt{DEF-07}), la app mostró ``Completada · 1 de 2 entregados · 1 con
intento fallido · Sigues ocupado: regresa al depósito y confirma tu llegada''.}{PASS}{JSON:
\texttt{evidencias/json/CP-077.json} y \texttt{movil-telefono.json}; captura del teléfono}{Ninguno}

\casoejec{CP-078}{Bloqueo de nueva salida antes de confirmar regreso (returnedAt)}{Con la ruta previa
\texttt{COMPLETED} y \texttt{returnedAt} nulo, asignar a DRIVER-A una nueva ruta de hoy respondió 409.
La segunda ruta quedó en \texttt{DRAFT} sin conductor y DRIVER-A conservó una sola reserva activa, sin
doble reserva.}{PASS}{JSON: \texttt{evidencias/json/CP-078.json}}{Ninguno}

\casoejec{CP-079}{Confirmación de regreso por titular libera y actualiza disponibilidad}{\texttt{POST
/routes/:id/return} de DRIVER-A respondió 200 con \texttt{returnedAt}, liberó la reserva (0 activas) y
dejó al repartidor disponible. Con otro repartidor con ausencia vigente hoy, el regreso también se
registró y liberó la reserva, pero su disponibilidad siguió en \texttt{false} por la ausencia. En el
teléfono, tras confirmar el regreso por API (\texttt{DEF-07}), la app mostró ``Visitas finalizadas.
Regreso al depósito confirmado'' y el repartidor quedó disponible.}{PASS}{JSON:
\texttt{evidencias/json/CP-079.json} y \texttt{movil-telefono.json}; captura del teléfono}{Ninguno}

\casoejec{CP-080}{Repetición de regreso a tienda sin duplicar ni liberar dos veces}{Repetir
\texttt{POST /routes/:id/return} respondió 200 con exactamente la misma respuesta que la primera vez
(idempotente). \texttt{returnedAt} y \texttt{returnedById} no cambiaron, no se creó ni liberó otra
reserva y la disponibilidad del repartidor siguió igual.}{PASS}{JSON y consulta a la base:
\texttt{evidencias/json/CP-080.json}}{Ninguno}

\casoejec{CP-081}{Intento de cierre o regreso por DRIVER-B rechazado sobre ruta ajena}{Con el JWT de
DRIVER-B, \texttt{POST /routes/:id/complete} y \texttt{POST /routes/:id/return} sobre la ruta de
DRIVER-A respondieron 403. El estado, \texttt{completedAt}, \texttt{returnedAt} y la reserva activa
de DRIVER-A quedaron exactamente iguales.}{PASS}{JSON: \texttt{evidencias/json/CP-081.json}}{Ninguno}

\clearpage
% CALIDAD - fichas individuales de ejecución (26 casos: CP-082..CP-107)
% GUÍA EQUIPO (José Mateo Auz): completar los parámetros de \casoejec{ID}{título}{resultado obtenido}{PASS|FAIL|BLOCKED|N/A}{evidencia o log}{DEF-XX}.
% - Estado: solo PASS, FAIL, BLOCKED o N/A.
% - Defecto: ID del defecto (ej. DEF-01) si falló, o Ninguno si pasó.
% - Resultado obtenido: describir lo observado al ejecutar la prueba.
% - Si existe evidencias/<ID>.png, la imagen se inserta automáticamente dentro de su ficha.
\par\medskip\noindent\textbf{Módulo Calidad --- Ruteo, incidencias, concurrencia y rendimiento (26 casos: CP-082 a CP-107).}\par\smallskip
Pruebas ejecutadas sobre planificación de rutas, gestión de incidencias y emergencias, concurrencia y mediciones de rendimiento (\texttt{100}/\texttt{500}/\texttt{2000} pedidos).

\casoejec{CP-082}{Agrupación horaria 10:\allowbreak{}00+10:\allowbreak{}20 y separación 15:\allowbreak{}00}{Con el motor vial OSRM real (distanceMethod \texttt{OSRM\_\allowbreak{}ROAD\_\allowbreak{}NETWORK}, margen 30 min) y tres pedidos del mismo día, \texttt{GET /route-proposals} como ADMIN devolvió HTTP 200 con 2 rutas: la ruta 1 agrupa P1 (10:00) y P2 (10:20, distante, 38,5 km) y la ruta 2 contiene solo P3 (15:00), aunque P3 queda a unos metros de P1. No se mezclaron por cercanía.}{PASS}{Log HTTP y JSON de propuesta: evidencias/json/CP-082.json}{Ninguno}
\casoejec{CP-083}{Ruta válida con 1 parada sin exigir mínimo de 4}{Con un único pedido elegible, \texttt{GET /route-proposals} devolvió HTTP 200 con 1 ruta de 1 parada (1,607 km por la red vial real, método \texttt{OSRM\_\allowbreak{}ROAD\_\allowbreak{}NETWORK}); no se exigió el mínimo de 4 pedidos ni hubo error de capacidad.}{PASS}{Log HTTP y JSON de propuesta: evidencias/json/CP-083.json}{Ninguno}
\casoejec{CP-084}{Lote de 4 pedidos compatibles forma ruta de 4 paradas}{Cuatro pedidos compatibles (10:00, 10:10, 10:15 y 10:25) se agruparon en 1 sola ruta de exactamente 4 paradas, en orden cronológico (A $\rightarrow$ B $\rightarrow$ C $\rightarrow$ D), con 17,617 km frente a 20,186 km de la línea base (12,73 \% de ahorro).}{PASS}{Log HTTP y JSON de propuesta: evidencias/json/CP-084.json}{Ninguno}
\casoejec{CP-085}{Lote de 5 pedidos dividido en máximo 4 paradas por ruta}{Cinco pedidos compatibles se dividieron en 2 rutas (1 + 4 paradas): ninguna ruta tiene 5 paradas, la suma es 5 y no hay pedidos duplicados ni excluidos (plannedOrders 5, excludedOrders 0).}{PASS}{Log HTTP y JSON de propuesta: evidencias/json/CP-085.json}{Ninguno}
\casoejec{CP-086}{Declarar emergencia pasa ruta a INTERRUPTED y bloquea}{\texttt{EMERGENCY\_\allowbreak{}DECLARE} sobre una ruta \texttt{IN\_\allowbreak{}PROGRESS} de 4 paradas (1 entregada, 1 fallida, 2 pendientes): la ruta pasó a \texttt{INTERRUPTED} sin \texttt{completedAt} ni \texttt{returnedAt}; se conservaron las paradas SUCCESSFUL y FAILED con sus intentos; las reservas se liberaron en el instante de la emergencia, los pedidos pendientes quedaron en PENDING con revisión administrativa y se creó un \texttt{RouteUnitRecovery} por cada unidad pendiente (4). Repetir la misma petición devolvió la misma auditoría; otro actor o cuerpo distinto dio 409; una segunda emergencia dio 409 \texttt{EMERGENCY\_\allowbreak{}ALREADY\_\allowbreak{}DECLARED}. DRIVER y CLIENT 403, anónimo 401 y \texttt{actorId} en el cuerpo 400.}{PASS}{Suite emergency-recovery (S1 y S7), log Jest}{Ninguno}
\casoejec{CP-087}{Retorno individual desbloquea solo unidad confirmada}{\texttt{UNIT\_\allowbreak{}RECOVERY\_\allowbreak{}CONFIRM} sobre la unidad 1 de un pedido de 2 unidades: la unidad 1 quedó recuperada (quién y cuándo) y la 2 siguió pendiente con el pedido bloqueado (\texttt{orderAllUnitsResolved} falso). Repetir la petición devolvió la misma respuesta; otro \texttt{requestId} sobre la unidad ya recuperada dio 409 \texttt{UNIT\_\allowbreak{}ALREADY\_\allowbreak{}RECOVERED}. \texttt{Product.stock}, comprobante, intento de pago y total del pedido no cambiaron, y CONTACT, PICKUP\_AGREE, reencolado y planificación del pedido bloqueado respondieron 409 \texttt{RECOVERY\_\allowbreak{}PENDING}.}{PASS}{Suite emergency-recovery (S2 y S5), log Jest}{Ninguno}
\casoejec{CP-088}{Cierre de emergencia con pendientes mantiene bloqueo 409}{\texttt{EMERGENCY\_\allowbreak{}CLOSE} con una unidad sin recuperar (2 unidades, 1 recuperada): HTTP 200 con \texttt{closedAt} y \texttt{closedById}; repetirlo con la misma clave devolvió lo mismo y con otra clave 409. El detalle conservó \texttt{totalUnits} 2, \texttt{recoveredUnits} 1 y \texttt{pendingUnits} 1 y la unidad pendiente siguió listada en \texttt{GET /dispatch/emergencies}. Gestionar el pedido afectado dio 409 \texttt{RECOVERY\_\allowbreak{}PENDING} hasta confirmar la segunda unidad.}{PASS}{Suite emergency-recovery (S3), log Jest}{Ninguno}
\casoejec{CP-089}{Cierre de emergencia no libera repartidor si returnedAt NULL}{Con la ruta \texttt{INTERRUPTED} y una unidad con \texttt{recoveredAt} NULL, \texttt{EMERGENCY\_\allowbreak{}CLOSE} respondió 200 pero la ruta siguió con \texttt{returnedAt} NULL, la \texttt{DriverReservation} con \texttt{releasedAt} NULL y el repartidor no disponible (\texttt{isAvailable} falso, \texttt{hasPersistedHolds} verdadero, ruta bloqueante listada). El cierre administrativo no libera al repartidor.}{PASS}{JSON de estado en BD: evidencias/json/CP-089.json}{Ninguno}
\casoejec{CP-090}{Liberación de repartidor por ADMIN respeta ausencias y 409}{Ruta A: mientras \texttt{returnedAt} es NULL el repartidor no está disponible y una nueva asignación dio 409; el retorno propio fijó \texttt{returnedAt} y \texttt{returnedById} y liberó la reserva (\texttt{releasedAt} igual a \texttt{returnedAt}); una liberación posterior del ADMIN dio 409 \texttt{DRIVER\_\allowbreak{}ALREADY\_\allowbreak{}RELEASED}. Ruta B con ausencia vigente (\#44): \texttt{EMERGENCY\_\allowbreak{}DRIVER\_\allowbreak{}RELEASE} del ADMIN liberó la reserva una sola vez (retorno por el ADMIN), el retorno del repartidor tras la liberación y una segunda liberación dieron 409.}{PASS}{Suite emergency-recovery (S4), log Jest}{Ninguno}
\casoejec{CP-091}{Declarar pérdida con stock \ensuremath{\geq} 1 prepara reemplazo y descuenta}{\texttt{LOSS\_\allowbreak{}DECLARE} con \texttt{Product.stock} = 2: reemplazo \texttt{PREPARED} sin esperar el retorno del original, stock 2 $\rightarrow$ 1 y un movimiento \texttt{REPLACEMENT\_\allowbreak{}OUT} (delta $-$1, \texttt{stockAfter} 1, actor ADMIN). Pedido, compras, cobros y comprobantes sin cambio y el seguimiento muestra el reemplazo preparado y el original perdido.}{PASS}{Suite loss-replacement, log Jest}{Ninguno}
\casoejec{CP-092}{Reemplazo mantiene mismo pedido y no genera nuevo cobro}{Sobre un pedido con 1 \texttt{PurchaseReceipt}, 1 intento de pago y saldo de tarjeta 500000 c, \texttt{LOSS\_\allowbreak{}DECLARE} con reemplazo preparado dejó los mismos conteos (1 pedido, 1 comprobante, 1 pago), el saldo sin cambio (500000 c) y el mismo identificador de pedido; el historial muestra la pérdida con su reemplazo \texttt{PREPARED}. No hubo nuevo cobro ni comprobante.}{PASS}{JSON antes/después en BD: evidencias/json/CP-092.json}{Ninguno}
\casoejec{CP-093}{Declarar pérdida con stock 0 crea PENDING sin stock negativo}{\texttt{LOSS\_\allowbreak{}DECLARE} con \texttt{Product.stock} = 0: pérdida creada con reemplazo \texttt{PENDING}, stock sigue en 0 (nunca negativo) y cero movimientos \texttt{REPLACEMENT\_\allowbreak{}OUT}; no se simuló ningún reemplazo preparado.}{PASS}{Suite loss-replacement, log Jest}{Ninguno}
\casoejec{CP-094}{Contacto y reprogramación tras reemplazo sin entrega hoy}{Tras una entrega fallida y la ruta terminada, \texttt{LOSS\_\allowbreak{}DECLARE} dejó el reemplazo \texttt{PREPARED} (stock 3 $\rightarrow$ 2). Luego \texttt{CONTACT} con acuerdo explícito y \texttt{RESCHEDULE} con fecha futura (2026-11-16 12:00 Guayaquil) respondieron 200 y quedaron auditados con el mismo actor (\texttt{LOSS\_\allowbreak{}DECLARE}, \texttt{CONTACT} y \texttt{RESCHEDULE}, todos APPLIED). El pedido quedó PENDING con la fecha acordada, sin ruta ni revisión pendiente, no se marcó entregado y no se creó otra parada (solo la fallida original).}{PASS}{JSON de estado y auditoría en BD: evidencias/json/CP-094.json}{Ninguno}
\casoejec{CP-095}{Preparación diferida al reponer stock y 409 NO\_FREE\_STOCK}{Con stock 0, \texttt{REPLACEMENT\_\allowbreak{}PREPARE} respondió 409 \texttt{NO\_\allowbreak{}FREE\_\allowbreak{}STOCK} sin escribir nada (stock 0, sin auditoría nueva, estado PENDING). Tras reponer el stock a 1 el reemplazo pasó a \texttt{PREPARED} y el stock a 0 con un movimiento; repetirlo dio 409 \texttt{REPLACEMENT\_\allowbreak{}ALREADY\_\allowbreak{}PREPARED} sin otro descuento.}{PASS}{Suite loss-replacement, log Jest}{Ninguno}
\casoejec{CP-096}{Recuperación tardía SELLABLE suma +1 sin doble entrega}{Ejemplo del issue: 2 libres $\rightarrow$ el reemplazo deja 1 $\rightarrow$ \texttt{ORIGINAL\_\allowbreak{}RECOVER} \texttt{SELLABLE} suma una unidad (stock 2) con movimientos \texttt{REPLACEMENT\_\allowbreak{}OUT} y \texttt{RECOVERY\_\allowbreak{}IN}; el pedido no cambia de estado y la recuperación no crea rutas ni otra entrega. Con el reemplazo aún pendiente, la recuperación vendible habilita preparar después el reemplazo con ese stock.}{PASS}{Suite loss-replacement, log Jest}{Ninguno}
\casoejec{CP-097}{Recuperación tardía DAMAGED registra estado sin sumar stock}{\texttt{ORIGINAL\_\allowbreak{}RECOVER} \texttt{DAMAGED}: HTTP 200 con condición DAMAGED y actor registrados, stock sin aumento (se mantiene en 1) y sin movimiento \texttt{RECOVERY\_\allowbreak{}IN}; reclasificarlo como vendible o repetir la acción dio 409 y el stock no cambió.}{PASS}{Suite loss-replacement, log Jest}{Ninguno}
\casoejec{CP-098}{Repetición de ORIGINAL\_RECOVER rechaza 409 sin doble suma}{Una segunda confirmación (otro administrador, nueva clave) dio 409 \texttt{ORIGINAL\_\allowbreak{}ALREADY\_\allowbreak{}RECOVERED} y el stock quedó igual (2); con cuatro confirmaciones simultáneas una respondió 200 y tres 409, con un solo movimiento \texttt{RECOVERY\_\allowbreak{}IN}.}{PASS}{Suite loss-replacement, log Jest}{Ninguno}
\casoejec{CP-099}{Cancelación LOSS\_CANCEL restituye stock y bloquea si no aplica}{\texttt{LOSS\_\allowbreak{}CANCEL} sobre una pérdida con reemplazo \texttt{PREPARED} devolvió una unidad al stock con un movimiento \texttt{REPLACEMENT\_\allowbreak{}CANCEL\_\allowbreak{}IN}, fue idempotente y permitió volver a declarar la misma unidad; sobre una pérdida \texttt{PENDING} (stock 0) no generó movimientos y el stock siguió en 0; reactivó la revisión administrativa cuando el pedido tenía un intento fallido ya despejado. Se rechazó con 409 cuando el original ya fue recuperado, por estado del pedido, por reclamo activo o por retiro acordado, y la base de datos impide anular con \texttt{recoveredAt}. Con anulaciones concurrentes solo una repuso la unidad y las demás dieron 409. El cruce con la emergencia y el retiro en tienda (\#51 y \#50) también pasó.}{PASS}{Suites loss-replacement y emergency-recovery (S6), log Jest}{Ninguno}
\casoejec{CP-100}{Concurrencia UNIT\_RECOVERY\_CONFIRM vs LOSS\_DECLARE en unidad}{En 4 carreras de \texttt{UNIT\_\allowbreak{}RECOVERY\_\allowbreak{}CONFIRM} contra \texttt{LOSS\_\allowbreak{}DECLARE} sobre la misma unidad con claves distintas, siempre hubo exactamente un 200 y un 409 específico: en la carrera 1 ganó la recuperación (la pérdida dio 409 \texttt{UNIT\_\allowbreak{}ALREADY\_\allowbreak{}RECOVERED}) y en las carreras 2 a 4 ganó la pérdida (la recuperación dio 409 \texttt{UNIT\_\allowbreak{}HAS\_\allowbreak{}ACTIVE\_\allowbreak{}LOSS}). El estado persistido nunca quedó doble: una unidad recuperada sin pérdida activa, o una pérdida activa sin recuperación.}{PASS}{JSON de las 4 carreras: evidencias/json/CP-100.json}{Ninguno}
\casoejec{CP-101}{Idempotencia de intervenciones con requestId UUID y conflicto}{La misma petición con el mismo \texttt{requestId} devolvió la misma respuesta y auditoría y descontó una sola vez; con otro cuerpo la misma clave dio 409. En despacho, la repetición exacta devolvió el mismo registro de \texttt{DispatchIntervention} y las cargas concurrentes con claves iguales y cuerpos distintos persistieron una sola decisión.}{PASS}{Suites loss-replacement y urgent-dispatch, log Jest}{Ninguno}
\casoejec{CP-102}{Rendimiento generación de rutas con volumen de 100 pedidos}{Con 100 pedidos del conjunto sintético (huella 703324b68f51) y el motor vial real, \texttt{GET /route-proposals} respondió en 7 de 7 repeticiones: p50 233 ms, p95 366 ms y máximo 366 ms, 0 errores, método \texttt{OSRM\_\allowbreak{}ROAD\_\allowbreak{}NETWORK} y 26 rutas propuestas (ahorro vial 61,65 \%). Medido con Apache JMeter 5.6.3 (Java 25), 1 usuario. Entorno: Intel Core i9-14900HX, 31,7 GB, Node 26.11.1, PostgreSQL 16 y OSRM v5.26.0 en la misma laptop. Sin umbral SLA.}{PASS}{Reporte JMeter: evidencias/perf/jmeter-limite-300-100.jtl y volumen-jmeter-limite-300.json}{Ninguno}
\casoejec{CP-103}{Rendimiento generación de rutas con volumen de 500 pedidos}{Con la configuración por defecto (límite de 300 puntos del motor vial) 500 pedidos no se pueden planificar: las 7 peticiones respondieron 422 \texttt{ROAD\_\allowbreak{}MATRIX\_\allowbreak{}TOO\_\allowbreak{}LARGE} (``505 puntos y el límite es 300'') en p50 27 ms. Subiendo \texttt{OSRM\_\allowbreak{}MAX\_\allowbreak{}POINTS} a 2100, 500 pedidos se planifican con 0 errores en 7 repeticiones: p50 3384 ms, p95 4235 ms, máximo 4235 ms. Mismo entorno y versión que en 100 pedidos; sin umbral SLA.}{PASS}{Reporte JMeter: evidencias/perf/jmeter-*-500.jtl y volumen-jmeter-*.json}{Ninguno}
\casoejec{CP-104}{Rendimiento generación de rutas con volumen de 2000 pedidos}{Primera ejecución: FAIL. Con el motor real y \texttt{OSRM\_\allowbreak{}MAX\_\allowbreak{}POINTS} elevado, 2000 pedidos dieron 502 \texttt{ROAD\_\allowbreak{}ENGINE\_\allowbreak{}BAD\_\allowbreak{}RESPONSE} en 7 de 7 repeticiones porque OSRM devuelve 8 distancias levemente negativas ($-0{,}2$ a $-0{,}6$ m) entre 4\,020\,025 pares y el backend descartaba toda la propuesta (DEF-02). Corregido ese defecto, 2 de 7 repeticiones dieron 503 \texttt{ROAD\_\allowbreak{}ENGINE\_\allowbreak{}UNAVAILABLE} por un cierre de conexión aislado sin reintento (DEF-04). Reejecución con ambas correcciones (PR \#76, ya en \texttt{main}: commit 9b8557c), JMeter 5.6.3, 7 repeticiones por volumen: 100 pedidos p50 219 ms (máx. 301), 500 pedidos p50 3773 ms (máx. 4058) y 2000 pedidos p50 52978 ms (p95 55804, máx. 55804), todas con 0 errores; con 2000 pedidos se planificaron los 2000 en 505 rutas. Con el límite por defecto (300 puntos) 2000 pedidos siguen dando 422 \texttt{ROAD\_\allowbreak{}MATRIX\_\allowbreak{}TOO\_\allowbreak{}LARGE}, una protección configurable. Sin umbral SLA.}{PASS}{Reporte comparativo JMeter: evidencias/perf/jmeter-sin-limite-corregido-*.jtl y volumen-jmeter-*.json}{DEF-02 y DEF-04 (corregidos)}
\casoejec{CP-105}{Concurrencia Shop última unidad: [201, 409 STOCK\_UNAVAILABLE]}{Con stock = 1 y dos solicitudes simultáneas de \texttt{POST /shop/checkout} de clientes y tarjetas distintos: exactamente una aprobada (201) y una rechazada (409 \texttt{STOCK\_\allowbreak{}UNAVAILABLE}), un solo pedido y comprobante, stock final 0 y sin débito ni pedido para la rechazada. En la matriz, 12 compradores simultáneos por 5 unidades: 5 aprobados, 7 con \texttt{STOCK\_\allowbreak{}UNAVAILABLE} y stock 0.}{PASS}{Suites shop-checkout-v12 y shop-concurrency-matrix-v12, log Jest}{Ninguno}
\casoejec{CP-106}{Concurrencia Shop saldo: [201, 422 INSUFFICIENT\_DEMO\_FUNDS]}{Tarjeta \texttt{DEMO\_\allowbreak{}LIMITED} con saldo 2000 c y dos compras simultáneas de 1725 c: una aprobada (201) y la otra rechazada (422 \texttt{INSUFFICIENT\_\allowbreak{}DEMO\_\allowbreak{}FUNDS}), saldo final 275 c, stock descontado una sola vez y nada cobrado ni creado para la rechazada. En la matriz, 10 compras simultáneas contra la misma tarjeta: solo una cabe y el saldo no queda negativo.}{PASS}{Suites shop-checkout-concurrency-v12 y shop-concurrency-matrix-v12, log Jest}{Ninguno}
\casoejec{CP-107}{Fallo pre-commit Shop con trigger en purchaseReceipt: rollback}{Con un disparador \texttt{BEFORE INSERT} controlado en \texttt{purchaseReceipt} (instalado solo en la base desechable y retirado al terminar), un checkout válido falló con HTTP 500 y deshizo todo: los conteos de pedido, ítem, intento de pago y comprobante, el stock y el saldo demo quedaron idénticos al estado previo a la petición.}{PASS}{Suite shop-checkout-v12 (Late Write Failure Rollback Probe), log Jest}{Ninguno}


% ===========================================================================
\clearpage
\setcounter{secnumdepth}{3}
\subsubsection{Registro de incidentes}
% ===========================================================================
% REPORTE DE DEFECTOS
% GUÍA EQUIPO (Calidad consolida):
% - Todo caso con estado FAIL genera una fila de defecto en esta tabla.
% - Cada integrante registra los suyos y Calidad consolida el reporte final.
% - Cada defecto se reporta también como Issue en el repositorio correspondiente.
% - Severidad ISO/IEC/IEEE 29119-3: S1 bloqueante, S2 alta, S3 media, S4 baja.
% ===========================================================================
Registro consolidado de defectos identificados durante la ejecución de las pruebas, clasificados
según la escala de severidad del estándar ISO/IEC/IEEE 29119-3: \textbf{S1} bloqueante
(seguridad, datos o ruteo crítico), \textbf{S2} alta, \textbf{S3} media y \textbf{S4} baja.

\begin{center}
\footnotesize
\setlength{\tabcolsep}{3pt}
\begin{longtable}{|P{1.2cm}|P{1.5cm}|P{3.6cm}|P{3.2cm}|P{0.9cm}|P{1.9cm}|P{1.6cm}|}
\hline
\textbf{ID} & \textbf{Caso} & \textbf{Descripción y pasos para reproducir} & \textbf{Esperado / obtenido} & \textbf{Sev.} & \textbf{Responsable} & \textbf{Estado} \\ \hline
\endhead
DEF-01 & & & & & & \\ \hline
DEF-02 & & & & & & \\ \hline
DEF-03 & & & & & & \\ \hline
DEF-04 & & & & & & \\ \hline
DEF-05 & & & & & & \\ \hline
DEF-06 & & & & & & \\ \hline
DEF-07 & & & & & & \\ \hline
DEF-08 & & & & & & \\ \hline
\end{longtable}
\end{center}


% ===========================================================================
\subsubsection{Verificación de resultados}
% GUÍA PARA EL EQUIPO (Calidad consolida):
% 1. Sumar PASS, FAIL, BLOCKED y N/A por cada módulo y calcular % de éxito = PASS / ejecutados.
% 2. Completar las cifras totales y la conclusión en el párrafo inferior antes de entregar el PDF.
Resumen consolidado de verificación frente a los criterios de aceptación del sistema. Un caso se
considera exitoso (\texttt{PASS}) cuando cumple íntegramente el criterio establecido sin efectos
parciales.

\begin{center}
\begin{tabularx}{\textwidth}{|l|c|c|c|c|c|L|}
\hline
\textbf{Módulo / Área} & \textbf{Casos} & \textbf{PASS} & \textbf{FAIL} & \textbf{BLOCKED} & \textbf{N/A} & \textbf{\% de éxito (PASS / ejecutados)} \\ \hline
Backend y reglas de negocio & 27 & 27 & 0 & 0 & 0 & 100\% (27 de 27; última ejecución) \\ \hline
Aplicación Web REGALÍA & 27 & 27 & 0 & 0 & 0 & 100\% (27 de 27) \\ \hline
Aplicación Móvil & 27 & 24 & 3 & 0 & 0 & 88,89\% (24 de 27) \\ \hline
Calidad y Rendimiento & 26 & 26 & 0 & 0 & 0 & 100\% (26 de 26) \\ \hline
\textbf{Total} & \textbf{107} & \textbf{104} & \textbf{3} & \textbf{0} & \textbf{0} & \textbf{97,20\% (104 de 107)} \\ \hline
\end{tabularx}
\end{center}

% GUÍA PARA EL EQUIPO: reemplazar los espacios subrayados con los conteos reales PASS/FAIL/BLOCKED/N/A, defectos S1..S4 y conclusión final.
De los 107 casos registrados, 104 obtuvieron \texttt{PASS} y 3 \texttt{FAIL}; no se registraron
casos \texttt{BLOCKED} ni \texttt{N/A}. Los tres fallos corresponden a \texttt{CP-067},
\texttt{CP-071} y \texttt{CP-073}. El reporte contiene nueve incidencias: cuatro abiertas,
tres corregidas y dos cerradas por aclaración de fichas. Por severidad se registran dos S2,
dos S3 y tres S4; las dos aclaraciones no tienen severidad de producto.\par
Las acciones de cierre y retorno de algunos escenarios móviles se realizaron mediante la API,
por lo que sus resultados no acreditan la disponibilidad de dichas acciones en la interfaz móvil.
Los defectos abiertos S2 asociados a la operación móvil limitan la aceptación integral del sistema.

\medskip\noindent\textbf{Mediciones de planificación.} Con \texttt{OSRM\_MAX\_POINTS=2100},
OSRM real y siete repeticiones por volumen, la mediana registrada fue de 219 ms para 100 pedidos,
3773 ms para 500 y 52978 ms para 2000; el máximo para 2000 fue 55804 ms. En la configuración
con límite de 300 puntos, los lotes de 500 y 2000 devolvieron HTTP 422. Estas mediciones describen
el entorno indicado y no se comparan con un umbral de aceptación no definido.

% ===========================================================================
\subsubsection{Entrega de evidencias}
% GUÍA PARA EL EQUIPO: como solo se entrega el PDF final, no se duplican capturas aquí ni se referencia una carpeta externa.
% Todas las capturas (evidencias/CP-XXX.png) y extractos de logs/reportes deben quedar embebidos en 2.7.4.
Las imágenes y extractos de evidencia se integran en las fichas correspondientes de la
sección~2.7.4. El presente documento consolida los resultados y las evidencias en un único PDF.

\clearpage
\subsection{Habilidades blandas empleadas en la práctica}
 $\boxtimes$ Liderazgo \\
 $\boxtimes$ Trabajo en equipo \\
 $\square$ Comunicación asertiva \\
 $\square$ La empatía \\
 $\square$ Pensamiento crítico \\
 $\square$ Flexibilidad \\
 $\boxtimes$ La resolución de conflictos \\
 $\square$ Adaptabilidad \\
 $\boxtimes$ Responsabilidad \\

\subsection{Conclusiones}
% GUÍA PARA EL EQUIPO: si se desea, agregar cifras exactas (% de éxito y latencias observadas en 100/500/2000 pedidos) tras llenar las tablas.
\begin{itemize}
    \item La última ejecución registrada de los 107 casos produjo 104 resultados satisfactorios y 3 fallidos, equivalentes al 97,20\% y al 2,80\%, respectivamente. Los fallos se concentraron en los escenarios móviles \texttt{CP-067}, \texttt{CP-071} y \texttt{CP-073}.
    \item Las pruebas registradas abarcaron comercio, acceso por roles, cobertura, planificación, incidencias, concurrencia y recuperación. La presencia de dos defectos móviles abiertos de severidad S2 impide concluir que la operación móvil cumple íntegramente los criterios evaluados.
    \item En el entorno de medición, el aumento del volumen de pedidos incrementó los tiempos observados de planificación. Para 2000 pedidos se obtuvo una mediana de 52978 ms con el límite ampliado; los resultados constituyen una línea base y no una evaluación frente a un SLA no establecido.
\end{itemize}

\subsection{Recomendaciones}
\begin{itemize}
    \item Integrar de forma continua la ejecución automatizada de las suites unitarias, de integración y extremo a extremo vinculadas a los identificadores \texttt{CP-001}--\texttt{CP-107}.
    \item Priorizar la corrección y revalidación de cualquier defecto de severidad \textbf{S1} o \textbf{S2} detectado durante las pruebas antes de su despliegue en producción.
    \item Mantener el monitoreo periódico de los tiempos de respuesta en la planificación de rutas para volúmenes crecientes (\texttt{100}, \texttt{500} y \texttt{2000} pedidos).
\end{itemize}

\subsection{Referencias bibliográficas}
\begin{enumerate}[label={[\arabic*]}]
    \item J. Mejía Trejo, \textit{Principios de aseguramiento de la calidad para el diseño de software}, 1.\textsuperscript{a} ed. Jalisco, México: AMIDI, 2024.
    \item T. Winters, \textit{Ingeniería de Software en Google}, 1.\textsuperscript{a} ed. MARCOMBO, 2022.
    \item J. M. Ortega Candel, \textit{Desarrollo seguro en ingeniería del software}, 1.\textsuperscript{a} ed. España: MARCOMBO, 2020.
    \item G. Pantaleo, \textit{Calidad en el desarrollo de software}, 2.\textsuperscript{a} ed. Buenos Aires: Alfaomega, 2016.
    \item B. Coronado García, \textit{Despliegue y puesta en funcionamiento de componentes software. UF1291}, 1.\textsuperscript{a} ed. Logroño, España: Tutor Formación, 2025.
\end{enumerate}

\subsection{Anexos}
Las imágenes y extractos de resultados de las pruebas de API, web, móvil y rendimiento se
encuentran integrados en las fichas de ejecución de la sección~2.7.4 del presente informe PDF.

\end{document}
